% =============================================================================
% CAPÍTULO 4 — DESARROLLO DE LA SOLUCIÓN
% =============================================================================
\chapter{Desarrollo de la Solución}
\label{ch:desarrollo}

% ─────────────────────────────────────────────────────────────────────────────
% NOTA PARA REDACCIÓN
%
%  Este capítulo cuenta la historia completa de CÓMO funciona la plataforma.
%  La narrativa sigue el mismo orden que un usuario recorrería el sistema:
%    1. Por qué existe (comparación AS-IS -> TO-BE)
%    2. Qué es y cómo está organizado (visión, arquitectura, modelo de datos)
%    3. El flujo principal: propuesta de asignatura (importar -> editar -> validar -> aprobar)
%    4. Las capacidades inteligentes (controles automáticos, LLM, RAG, MCP)
%    5. Los módulos institucionales (plan de estudios, docentes, evaluaciones, acreditación)
%    6. La infraestructura (API, frontend, seguridad)
%    7. Las decisiones técnicas clave
% ─────────────────────────────────────────────────────────────────────────────

El presente capítulo describe el desarrollo de \textsc{macau}: desde las
decisiones de arquitectura hasta la implementación de cada módulo funcional.
El sistema combina gestión documental estructurada, validación normativa
automatizada, asistencia de inteligencia artificial y trazabilidad de
procesos en una plataforma diseñada para el contexto de la gestión
académica universitaria.

La exposición parte del diagnóstico presentado en el capítulo anterior y
lo confronta con el modelo propuesto, mostrando de qué manera cada punto
de fricción del proceso actual fue convertido en un mecanismo sistemático
y trazable.
A partir de ahí se describe la arquitectura de la plataforma, el modelo
de datos que la sostiene y el ciclo completo de vida de una propuesta de
asignatura: desde su importación hasta su aprobación, pasando por la edición
asistida y la validación normativa.

Las secciones centrales del capítulo se dedican a las capacidades
inteligentes del sistema: los controles automáticos de consistencia,
la asistencia LLM para la redacción, el módulo de búsqueda semántica
mediante recuperación aumentada por generación (RAG) y la capa de
agentes basada en el \emph{Model Context Protocol}.
Se presentan luego los módulos institucionales que completan la plataforma
---plan de estudios, competencias, docentes, instrumentos evaluativos,
acreditación y notificaciones--- y se finaliza con la descripción de la
API REST, la interfaz de usuario, las medidas de seguridad y las
decisiones de diseño que condicionaron la arquitectura.

%% =============================================================================
\section{Del proceso manual a la plataforma: contraste \textit{AS-IS} / \textit{TO-BE}}
\label{sec:asis_tobe}
%% =============================================================================

El diagnóstico presentado en la Sección~\ref{sec:proceso_asis} identificó seis dimensiones de
problema en el modelo de gestión académica vigente: ausencia de visibilidad sobre el estado
de las propuestas; coordinación dependiente de hilos de correo electrónico no estructurados;
falta de trazabilidad en el historial de versiones y aprobaciones; heterogeneidad en la
aplicación de los criterios de revisión; carga operativa acumulada en el director de carrera;
y baja escalabilidad ante el volumen de 45 a 55 asignaturas por ciclo lectivo
(véase Tabla~\ref{tab:problemas_asis}).
La plataforma desarrollada en este trabajo ---denominada \textsc{macau}
(\textit{MultiAgente de Apoyo a la Calidad Académica Universitaria})---
transforma cada uno de esos puntos de fricción
en un mecanismo sistemático, trazable y asistido por inteligencia artificial.

Para organizar este análisis se utiliza la notación \textbf{\textit{AS-IS} / \textit{TO-BE}},
habitual en el rediseño de procesos de negocio:
el modelo \textit{AS-IS} describe el proceso \emph{tal como funciona actualmente},
con sus herramientas, actores y fricciones reales;
el modelo \textit{TO-BE} describe cómo ese mismo proceso queda estructurado
\emph{una vez incorporada la plataforma}.
La comparación no evalúa tecnología por tecnología, sino dimensión por dimensión
del problema operativo identificado en el diagnóstico.

La Tabla~\ref{tab:asis_tobe} sintetiza la transformación en cada dimensión.
No se trata de digitalizar el proceso existente, sino de rediseñarlo: los instrumentos
de coordinación y control pasan de ser externos al proceso ---el correo electrónico,
la planilla Excel--- a ser intrínsecos a él, generados y registrados por el propio sistema
en cada acción de cada actor.

\setlength{\extrarowheight}{4pt}
\arrayrulecolor{black}
\begin{longtable}{p{3cm} p{5cm} p{7cm}}
\caption{Comparación del proceso de gestión académica antes y después de la plataforma.}
\label{tab:asis_tobe} \\
\toprule
\textbf{Dimensión} & \textbf{Proceso manual (\textit{AS-IS})} & \textbf{Con la plataforma (\textit{TO-BE})} \\
\midrule
\endfirsthead
\multicolumn{3}{r}{\small\textit{(continúa\ldots)}} \\[2pt]
\toprule
\textbf{Dimensión} & \textbf{Proceso manual (\textit{AS-IS})} & \textbf{Con la plataforma (\textit{TO-BE})} \\
\midrule
\endhead
\bottomrule
\endfoot
Entrega de propuesta
  & Docente envía DOCX por email o deposita en Google Drive;
    canal sin consistencia, versiones duplicadas posibles.
  & Dirección o Secretaría importa DOCX o enlaza Google Doc directamente
    en la plataforma; el sistema registra fecha, versión y autor. \\
\arrayrulecolor{gray!35}\hline\arrayrulecolor{black}
Seguimiento del estado
  & Director actualiza manualmente una planilla Excel;
    el docente no tiene visibilidad sin consultarlo explícitamente.
  & Dashboard en tiempo real por carrera: cada propuesta
    muestra su estado actual (Borrador / En revisión / Observada /
    Aprobada / Archivada) para todos los actores involucrados. \\
\arrayrulecolor{gray!35}\hline\arrayrulecolor{black}
Revisión de la CC
  & Cada miembro aplica el checklist según su criterio personal;
    resultados no homogéneos entre asignaturas ni entre ciclos.
  & Controles normativos autom\'{a}ticos (derivados de los criterios
    de acreditaci\'{o}n CONEAU) ejecutados por el sistema; reporte
    estructurado de conformidad por campo y por asignatura. \\
\arrayrulecolor{gray!35}\hline\arrayrulecolor{black}
Devolución de observaciones
  & Correo electrónico libre, con o sin copia al director;
    sin formato estandarizado ni registro centralizado.
  & Las observaciones se registran en la plataforma asociadas
    a la propuesta; el docente las recibe con una notificación
    automática generada por el Director o Miembro de la Comisión Curricular con asistencia LLM. \\
\arrayrulecolor{gray!35}\hline\arrayrulecolor{black}
Historial y versionado
  & Ninguno: el reenvío de una propuesta corregida sobreescribe
    silenciosamente la versión anterior en Drive.
  & Log de auditoría inmutable: cada cambio de estado,
    edición de campo o acción de agente queda registrado
    con actor, timestamp y hash del contenido. \\
\arrayrulecolor{gray!35}\hline\arrayrulecolor{black}
Instrumentos evaluativos
  & TPs, parciales y finales se comparten por email o Drive
    ad hoc; el director no tiene visibilidad sin consultar
    individualmente a cada docente.
  & Módulo centralizado de instrumentos por asignatura
    y carrera, vinculado a carpetas de Google Drive;
    dashboard de cumplimiento para el director. \\
\arrayrulecolor{gray!35}\hline\arrayrulecolor{black}
Evidencia para CONEAU
  & Esfuerzo concentrado al inicio del proceso de acreditación;
    recuperación de documentos dispersa e incompleta.
  & Registro continuo de evidencias con clasificación
    automática por dimensión CONEAU, versionado
    e historial de auditoría disponible en todo momento. \\
\end{longtable}
\arrayrulecolor{black}

La Figura~\ref{fig:flujo_tobe} ilustra el flujo \textit{TO-BE} del ciclo de gestión
de propuestas de asignatura.
Cabe señalar que este diagrama cubre exclusivamente ese proceso;
los flujos de instrumentos evaluativos, notificaciones institucionales y seguimiento
del plan de acreditación CONEAU son procesos diferenciados que se describen
en las secciones correspondientes de este capítulo.

Respecto al modelo \textit{AS-IS} (Figura~\ref{fig:flujo_asis}), el diagrama
incorpora tres aspectos que no existían en el proceso manual:
(i)~las dos vías de incorporación de una propuesta ---el director puede crearla
desde cero definiendo el encabezado institucional (nombre de la asignatura, equipo
docente, carga horaria, competencias) y el docente la completa desde la plataforma;
o bien la dirección o secretaría puede importar un DOCX existente o vincular
directamente un Google Doc---;
(ii)~el paso de edición con asistencia del agente LLM, que apoya al docente
con corrección ortográfica y gramatical, reformulación de secciones y generación
asistida de contenido desde el contexto de la asignatura; y
(iii)~la formalización del bucle de revisión: cada iteración tiene un estado
explícito, queda registrada en el log de auditoría y está asistida inteligentemente
tanto en la corrección del docente como en el proceso de validación de la CC.

\begin{figure}[H]
\centering
\scalebox{0.85}{%
\begin{tikzpicture}[
  node distance=0.5cm and 1.6cm,
  %% Estilo base (idéntico al de fig:llm_ciclo_vida): draw=#1!70, fill=#1!6
  mbox/.style={rectangle, rounded corners=6pt,
               draw=#1!70, fill=#1!6,
               text width=6.4cm, align=left, inner sep=8pt,
               font=\small},
  sbox/.style={rectangle, rounded corners=6pt,
               draw=#1!70, fill=#1!6,
               text width=3.4cm, align=center, inner sep=7pt,
               font=\small},
  dec/.style={diamond, draw=Micolor1!70, fill=Micolor1!6,
              font=\small\bfseries, text width=2.0cm, align=center,
              aspect=2.4, inner sep=2pt},
  arr/.style={-{Stealth[length=5pt]}, very thick, color=gray!50},
  lbl/.style={font=\scriptsize\bfseries, fill=white, inner sep=1pt}
]

%% ── Columna principal ────────────────────────────────────────────────────
\node[mbox=Micolor1] (entrada) {%
  \textbf{Incorporaci\'on de la propuesta}\\[3pt]
  \textbullet\ Director crea encabezado: asignatura, equipo, carga~hor., competencias\\[1pt]
  \textbullet\ Direcci\'on/Secr.\ importa DOCX o vincula Google Doc existente%
};

\node[mbox=teal, below=of entrada] (edita) {%
  \textbf{Edici\'on con asistencia LLM}\\[2pt]
  Redacci\'on \textbullet\ gram\'atica \textbullet\ reformulaci\'on \textbullet\ generaci\'on asistida%
};

\node[mbox=teal, below=of edita] (ctrl) {%
  \textbf{Controles normativos autom\'aticos}\\[2pt]
  Verificaci\'on campo a campo (criterios de acreditaci\'on CONEAU) \textbullet\ reporte de conformidad%
};

\node[dec, below=of ctrl] (ok1) {¿Cumple controles?};

\node[mbox=gray, below=of ok1] (envrev) {%
  \textbf{Propuesta enviada a revisi\'on CC}\\[2pt]
  Estado: \textit{En revisi\'on} \textbullet\ notificaci\'on autom\'atica%
};

\node[dec, below=of envrev] (ok2) {¿CC valida?};

\node[mbox=Micolor1, below=of ok2] (aprov) {%
  \textbf{Director aprueba}\\[2pt]
  Estado: \textit{Aprobada} \textbullet\ log de auditor\'ia inmutable%
};

%% ── Columna derecha — bucle de corrección ────────────────────────────────
\node[sbox=gray, right=1.6cm of ok1] (corr) {%
  \textbf{Docente corrige}\\[3pt]
  con asistente LLM\\[2pt]
  {\scriptsize sugerencias \textbullet\ reformulaci\'on}%
};

\node[sbox=teal, right=1.6cm of ok2] (obs) {%
  \textbf{CC registra\\observaciones}\\[2pt]
  {\scriptsize estructuradas\\en la plataforma}%
};

%% ── Flechas flujo principal ──────────────────────────────────────────────
\draw[arr] (entrada) -- (edita);
\draw[arr] (edita)   -- (ctrl);
\draw[arr] (ctrl)    -- (ok1);
\draw[arr] (ok1)     -- node[lbl,left]{S\'i} (envrev);
\draw[arr] (envrev)  -- (ok2);
\draw[arr] (ok2)     -- node[lbl,left]{S\'i} (aprov);

%% ── Bucle de corrección ──────────────────────────────────────────────────
\draw[arr] (ok1)       -- node[lbl,above]{No} (corr);
\draw[arr] (ok2)       -- node[lbl,above]{No} (obs);
\draw[arr] (obs.north) -- (corr.south);
\draw[arr] (corr.west) -- ++(-0.5,0) |- node[lbl,right,pos=0.25]{nueva versi\'on} (ctrl.east);

\end{tikzpicture}%
}% end scalebox
\caption{Flujo \textit{TO-BE} del ciclo de propuestas de asignatura en \textsc{macau}.
  A diferencia del modelo \textit{AS-IS} (Figura~\ref{fig:flujo_asis}), la plataforma
  incorpora dos v\'ias de entrada y asistencia LLM en la edici\'on y correcci\'on.
  \textit{Fuente: elaboraci\'on propia.}}
\label{fig:flujo_tobe}
\end{figure}

%% =============================================================================
\section{Visión general de la solución}
\label{sec:vision_general}
%% =============================================================================

\textsc{macau} es una plataforma web de gestión académico-administrativa asistida por inteligencia
artificial, orientada a centralizar, trazar y automatizar los procesos institucionales
vinculados a la calidad académica y a la acreditación universitaria.
Su alcance abarca desde el ciclo de vida de las propuestas de asignatura ---incluyendo
su elaboración, revisión, validación normativa y aprobación--- hasta la administración
de los trabajos prácticos e instrumentos de evaluación, el seguimiento del plan de
estudios y del catálogo de competencias, la documentación continua de evidencias para
la acreditación CONEAU, y la gestión de notificaciones institucionales con asistencia LLM.
No se trata de una herramienta de digitalización parcial, sino de un sistema que cubre
el ciclo documental completo de una unidad académica de ingeniería, desde el inicio
del ciclo lectivo hasta la presentación ante los organismos de acreditación.

La plataforma nació como respuesta al diagnóstico presentado en la
Sección~\ref{sec:proceso_asis}: un proceso disperso en correos electrónicos,
planillas Excel y versiones de documentos sin control de cambios que demanda
un esfuerzo operativo desproporcionado y deja escasa evidencia sistematizada.
Fue dise\~{n}ada, implementada y validada en el contexto de la
\textbf{Universidad Nacional de Chilecito (UNdeC)}, tomando como referencia normativa
la Ordenanza N\textdegree{}~75/23 de la CONEAU~\cite{ordenanza75}, que regula los
procedimientos para la acreditaci\'on de carreras de grado en funcionamiento,
y la reglamentaci\'on interna de la universidad sobre la presentaci\'on de programas
anal\'iticos ante los organismos competentes.

\subsection{Actores del sistema}

\textsc{macau} está concebido para cinco perfiles de usuario, que corresponden a los
actores identificados en la Sección~\ref{sec:proceso_asis}:

\begin{description}
  \item[Administrador] Gestiona la configuración global de la plataforma: crea y
    administra carreras, docentes y usuarios; configura integraciones externas
    (Google Workspace, OpenAI); y tiene acceso irrestricto a todos los módulos
    con fines de mantenimiento y auditoría del sistema.

  \item[Director de carrera] Supervisa el flujo completo del ciclo lectivo y el estado de
    todas las propuestas de su carrera mediante un \textit{dashboard} centralizado.
    Ejecuta o delega la revisión normativa, aprueba propuestas, gestiona la Comisión
    Curricular, envía notificaciones, configura controles inteligentes y tiene visibilidad agregada sobre
    el cumplimiento del programa académico (instrumentos entregados vs.\ pendientes), además de gestionar el proceso de acreditación generando y controlando el plan de trabajo.

  \item[Secretario de carrera] Apoya la gestión documental de la carrera: exporta
    propuestas en el formato oficial, administra el módulo de instrumentos evaluativos,
    coordina el registro de evidencias de acreditación CONEAU y hace seguimiento
    del plan de trabajo de la comisión de autoevaluación.

  \item[Docente] Elabora y entrega la propuesta de asignatura, asistido por el agente LLM
    para redacción, reformulación y verificación de cumplimiento normativo.
    Registra y sube los instrumentos de evaluación (trabajos prácticos, parciales y finales)
    de las asignaturas a su cargo.

  \item[Miembro de la Comisión Curricular (CC)] Accede a las propuestas en estado
    ``En revisión'', ejecuta los controles normativos automáticos como así también los controles inteligentes asistidos por LLM sobre los controles definidos por la carrera, agrega observaciones
    estructuradas y valida las propuestas para su aprobación final.
\end{description}

\subsection{Módulos funcionales}

\textsc{macau} integra siete módulos funcionales que cubren de extremo a extremo
los procesos de la unidad académica (véase la Tabla~\ref{tab:modulos_macau}).

\setlength{\extrarowheight}{4pt}
\arrayrulecolor{black}
\begin{longtable}{p{3.6cm} p{10.8cm}}
\caption{Módulos funcionales de \textsc{macau} y su correspondencia con los objetivos del sistema.}
\label{tab:modulos_macau} \\
\toprule
\textbf{Módulo} & \textbf{Capacidades principales} \\
\midrule
\endfirsthead
\multicolumn{2}{r}{\small\textit{(continúa\ldots)}} \\[2pt]
\toprule
\textbf{Módulo} & \textbf{Capacidades principales} \\
\midrule
\endhead
\bottomrule
\endfoot
Propuestas de asignatura
  & Creación, importación desde DOCX o Google Docs, edición por estados formales
    (Borrador / En revisión / Observada / Aprobada / Archivada), exportación
    en DOCX, PDF, XML y JSON, historial de versiones e historial de auditoría inmutable. \\
\arrayrulecolor{gray!35}\hline\arrayrulecolor{black}
Asistencia inteligente (LLM)
  & Generación asistida de secciones a partir del contexto de la propuesta,
    reformulación de texto con adecuación normativa, sugerencias de mejora campo
    a campo, controles de cumplimiento configurables mediante \textit{prompts} en
    lenguaje natural, y búsqueda semántica enriquecida con RAG sobre el corpus
    de propuestas y normativa vigente. \\
\arrayrulecolor{gray!35}\hline\arrayrulecolor{black}
Instrumentos de evaluación
  & Registro centralizado de trabajos prácticos, parciales, recuperatorios y finales
    por asignatura y carrera; seguimiento del estado de entrega por docente;
    vinculación a carpetas de Google Drive; \textit{dashboard} de cumplimiento
    del programa para el director; notificación automática a docentes con
    entregas pendientes. \\
\arrayrulecolor{gray!35}\hline\arrayrulecolor{black}
Plan de estudios, competencias y Matríz de Tributación
  & Gestión de la estructura jerárquica del plan (año $\to$ cuatrimestre $\to$ asignatura)
    con correlatividades; catálogo de competencias
    genéricas y específicas alineado con estándares CONFEDI/CONEAU; vinculación
    de competencias a propuestas y visualización de cobertura a nivel carrera; gestión de la Matríz de Tributación. \\
\arrayrulecolor{gray!35}\hline\arrayrulecolor{black}
Docentes y carreras
  & Gestión del catálogo institucional de docentes y carreras; asignación de roles
    diferenciados por carrera (Directivo / Docente / Secretario / Miembro\,CC);
    control de acceso granular por recurso y por estado del documento, con
    validación en cada operación de la API. \\
\arrayrulecolor{gray!35}\hline\arrayrulecolor{black}
Notificaciones institucionales
  & Envío de notificaciones por correo electrónico desde la cuenta institucional
    real vía Gmail OAuth~2.0: alertas sobre propuestas incompletas, instrumentos
    pendientes y vencimientos del plan de trabajo; generación del cuerpo del mail
    con asistencia LLM incluyendo contexto específico de asignatura y docente. \\
\arrayrulecolor{gray!35}\hline\arrayrulecolor{black}
Acreditación CONEAU
  & Registro, versionado y auditoría de evidencias documentales clasificadas por
    dimensión CONEAU; ingesta masiva desde Google Drive con clasificación automática
    por LLM; plan de trabajo de la comisión de autoevaluación con seguimiento de
    actividades y tareas; configuración del ciclo y estándar de acreditación vigente. \\
\end{longtable}
\arrayrulecolor{black}

La integración bidireccional con \textbf{Google Workspace} ---Google Docs para
la edición colaborativa de propuestas y Google Drive para el almacenamiento
documental--- atraviesa varios de estos módulos: el docente puede editar su propuesta
indistintamente desde la plataforma o desde el documento compartido, y el sistema
sincroniza los cambios campo a campo en ambas direcciones.
Esta integración elimina la duplicidad de herramientas que caracterizaba el proceso AS-IS
y garantiza coherencia permanente entre la versión colaborativa y el registro en la base
de datos institucional.

\subsection{Pilares tecnológicos}

La arquitectura de \textsc{macau} descansa en cuatro pilares tecnológicos
que se articulan entre sí (véase la Figura~\ref{fig:pilares_macau}):

\begin{description}
  \item[Sistema multi-agente (MAS).] Agentes especializados colaboran para
    cubrir funciones diferenciadas: extracción de documentos, asistencia de
    redacción, búsqueda semántica y notificación.
    La coordinación entre agentes se realiza desde la capa de negocio del backend.

  \item[Modelo de lenguaje grande (LLM).] La familia GPT-4o de OpenAI provee
    capacidade de generación, reformulación y evaluación en lenguaje natural.
    El contexto institucional ---carrera, perfil de egreso, normativa vigente---
    es inyectado en el \textit{system prompt} de cada interacción para mantener
    coherencia con el dominio.

  \item[Generación aumentada por recuperación (RAG).] Un índice vectorial
    (\textit{HNSW} sobre \texttt{hnswlib}) permite recuperar fragmentos
    relevantes de propuestas anteriores y documentos normativos antes de
    responder una consulta, reduciendo alucinaciones y enriqueciendo
    las respuestas con contexto específico del dominio.

  \item[Model Context Protocol (MCP).] Los agentes interactúan con los recursos
    del sistema a través de servidores MCP, que exponen herramientas y recursos
    normalizados. Este protocolo garantiza trazabilidad nativa: cada invocación
    de una herramienta queda registrada como evento auditable con identificador
    de agente, usuario, acción y \textit{timestamp} UTC.
\end{description}

\begin{figure}[H]
\centering
\scalebox{0.85}{%
\begin{tikzpicture}[
  pilar/.style={
    rectangle, rounded corners=6pt,
    draw=#1!70, fill=#1!6,
    font=\small\bfseries,
    text width=3.2cm, align=center,
    inner sep=8pt, minimum height=1.6cm
  },
  center/.style={
    rectangle, rounded corners=6pt,
    draw=Micolor1, line width=1.5pt,
    fill=Micolor1, text=white,
    font=\normalsize\bfseries,
    text width=2.6cm, align=center,
    minimum height=1.6cm
  },
  connarr/.style={
    -{Stealth[length=5pt]}, very thick, color=gray!50
  }
]
\node[center] (macau) {\textsc{macau}};
\node[pilar=Micolor1, above left=1.1cm and 2.0cm of macau]  (mas)  {MAS\\{\footnotesize Agentes especializados}};
\node[pilar=teal,     above right=1.1cm and 2.0cm of macau] (llm)  {LLM\\{\footnotesize GPT · asistencia}};
\node[pilar=teal,     below left=1.1cm and 2.0cm of macau]  (rag)  {RAG\\{\footnotesize Búsqueda semántica}};
\node[pilar=gray,     below right=1.1cm and 2.0cm of macau] (mcp)  {MCP\\{\footnotesize Trazabilidad · herramientas}};

\draw[connarr] (mas.east)  -- (macau.north west);
\draw[connarr] (llm.west)  -- (macau.north east);
\draw[connarr] (rag.east)  -- (macau.south west);
\draw[connarr] (mcp.west)  -- (macau.south east);
\end{tikzpicture}
}
\caption{Los cuatro pilares tecnológicos de \textsc{macau}.
  El sistema multi-agente (MAS), el modelo de lenguaje grande (LLM),
  la generación aumentada por recuperación (RAG) y el protocolo MCP
  se articulan en torno al núcleo de la plataforma.
  \textit{Fuente: elaboración propia.}}
\label{fig:pilares_macau}
\end{figure}

\subsection{Organización del capítulo}

Las secciones siguientes desarrollan en detalle cada uno de estos componentes.
La Sección~\ref{sec:arquitectura} describe la arquitectura en capas y el entorno
de despliegue. Las Secciones~\ref{sec:modelo_datos} a~\ref{sec:acreditacion}
cubren el modelo de datos y los módulos funcionales de la plataforma.
Las Secciones~\ref{sec:backend} y~\ref{sec:frontend} describen la implementación
del \textit{backend} REST y la interfaz web, respectivamente.
La Sección~\ref{sec:seguridad} analiza los mecanismos de seguridad y auditoría,
y la Sección~\ref{sec:decisiones_disenio} documenta las decisiones de diseño
más relevantes con sus justificaciones y compromisos.

%% =============================================================================
\section{Arquitectura general del sistema}
\label{sec:arquitectura}
%% =============================================================================

\textsc{macau} se estructura en cuatro capas lógicas que se comunican
mediante interfaces bien definidas (véase la
Figura~\ref{fig:arquitectura_capas}).
Esta separación permite reemplazar cualquier capa —por ejemplo, el cliente
web por una interfaz de línea de comandos— sin modificar el núcleo del
sistema.

\subsection{Capas de la arquitectura}
\label{subsec:capas}

\begin{description}

  \item[\textbf{Capa de presentación.}]
  La interfaz de usuario es una aplicación de página única
  (\textit{single-page application}, SPA) construida con
  \textbf{React} y \textbf{Vite}.
  El enrutamiento entre vistas se gestiona mediante renderizado condicional
  en el componente raíz \texttt{App.jsx}, sin dependencia de un enrutador
  externo.
  Los componentes consumen la API REST del \textit{backend} a través de la
  función nativa \texttt{fetch}, y el servidor de desarrollo de Vite escucha
  en el puerto~5173.

  \item[\textbf{Capa de lógica de negocio.}]
  El \textit{backend} es una aplicación \textbf{FastAPI} servida por
  \textbf{Uvicorn}, expuesta en el puerto~8011.
  La lógica de negocio se centraliza en \texttt{app/main.py}, que define los \textit{endpoints} REST organizados por
  recurso funcional: propuestas, docentes, carreras, controles inteligentes,
  exportación e importación.
  El middleware \texttt{CORSMiddleware} autoriza peticiones originadas en el
  cliente web.
  La autenticación se implementa con tokens \textbf{JWT} —biblioteca
  \texttt{PyJWT}, contraseñas protegidas con \texttt{bcrypt}—
  emitidos al iniciar sesión y verificados mediante inyección de dependencias
  en cada \textit{endpoint} protegido.

  \item[\textbf{Capa de agentes e integración.}]
  La plataforma incorpora múltiples componentes especializados que operan de
  forma autónoma o semi-autónoma sobre los datos del sistema.
  Formalmente, el directorio \texttt{agents/} aloja el módulo
  \texttt{extract.py}; sin embargo, el conjunto de responsabilidades
  agentivas es más amplio e incluye:

  \begin{itemize}[leftmargin=1.5em, itemsep=2pt]
    \item \textbf{Extracción de propuestas (DOCX/GDoc):}
          \texttt{agents/extract.py} en combinación con
          \texttt{app/docx\_import.py} analizan el documento subido e
          identifican automáticamente todos los campos curriculares
          (contenidos mínimos, resultados de aprendizaje, unidades
          temáticas, equipo docente, etc.) para poblar la base de datos.

    \item \textbf{Evaluación inteligente con LLM:}
          el módulo de controles inteligentes (\texttt{evaluate\_control\-\_with\_llm},
          \texttt{run\_intelligent\_controls}) aplica rúbricas configurables
          sobre cada propuesta consultando la API de OpenAI y devuelve
          retroalimentación estructurada con puntaje y observaciones.

    \item \textbf{Asistencia y reformulación LLM:}
          los \textit{endpoints} \texttt{ai\_generate}, \texttt{ai\_reformulate}
          y \texttt{suggest\_for\_proposal} permiten generar texto acad\'emico,
          reformular secciones y sugerir mejoras a partir de un contexto
          proporcionado por el usuario.

    \item \textbf{Búsqueda semántica (RAG):}
          \texttt{semantic\_search} localiza fragmentos normativos relevantes
          consultando el índice vectorial hnswlib y enriqueciéndolos con una
          respuesta generada por el LLM.

    \item \textbf{Gestor de instrumentos de evaluación:}
          \texttt{upload\_instruments} e \texttt{instruments\_notify}
          gestionan la carga de archivos a Google Drive y el envío
          automático de notificaciones por correo electrónico a los docentes
          correspondientes.

    \item \textbf{Sincronización con Google Workspace:}
          \texttt{sync\_proposal\_gdoc}, \texttt{import\_proposal\_gdoc\_url}
          e \texttt{ingest\_accreditation\_evidences} mantienen
          sincronizados los documentos entre la base de datos local y
          Google Drive / Docs, y realizan la ingesta de evidencias de
          acreditación desde carpetas compartidas.
  \end{itemize}

  \item[\textbf{Capa de persistencia.}]
  La base de datos relacional es \textbf{SQLite}, gestionada mediante el ORM
  \textbf{SQLAlchemy} con el patrón \textit{session-per-request}.
  Los modelos principales definidos en \texttt{app/models.py} son
  \texttt{Proposal} (propuesta curricular con todas sus secciones
  académicas), \texttt{Teacher} (docente), \texttt{Career} (carrera) y las
  tablas de asociación que registran la asignación de docentes y
  competencias a cada propuesta.
  Para la búsqueda semántica de fragmentos normativos, el sistema emplea
  \textbf{hnswlib}~\cite{malkov2018efficient}, una implementación
  eficiente de grafos jerárquicos de mundo pequeño
  (\textit{Hierarchical Navigable Small World}, HNSW) que permite consultas
  de vecinos más próximos sobre los vectores de incrustación generados por
  la API de OpenAI.

\end{description}

\begin{figure}[H]
\centering

\begin{tikzpicture}[
  layer/.style={
    rectangle, rounded corners=7pt,
    draw=#1!70, fill=#1!6,
    font=\small\bfseries,
    text width=13.0cm, align=center,
    inner sep=10pt, minimum height=1.7cm
  },
  %% agentes sin LLM directo
  subnode/.style={
    rectangle, rounded corners=5pt,
    draw=#1!60, fill=#1!10,
    font=\scriptsize,
    text width=2.1cm, align=center,
    inner sep=4pt, minimum height=1.0cm
  },
  extbox/.style={
    rectangle, rounded corners=6pt,
    draw=gray!50, fill=gray!5, dashed,
    font=\small\bfseries,
    text width=3.2cm, align=center,
    inner sep=8pt, minimum height=1.5cm
  },
  subcap/.style={
    rectangle, rounded corners=4pt,
    draw=gray!45, fill=gray!9,
    font=\scriptsize,
    text width=1.3cm, align=center,
    inner sep=3pt, minimum height=0.9cm
  },
  %% Navegador — violeta
  client/.style={
    rectangle, rounded corners=6pt,
    draw=violet!65, fill=violet!8,
    font=\small\bfseries,
    text width=2.2cm, align=center,
    inner sep=7pt, minimum height=1.4cm
  },
  arr/.style={
    {Stealth[length=5pt]}-{Stealth[length=5pt]}, thick, color=gray!60
  },
  %% flechas con dependencias externas — bidireccionales
  arrdash/.style={
    {Stealth[length=4pt]}-{Stealth[length=4pt]}, thick, dashed, color=gray!55
  }
]

%% ── Cliente (navegador) — arriba, centrado ─────────────────────────────────
\node[client] (client) at (0, 19.8)
  {Navegador\\[-1pt]{\normalfont\scriptsize(usuario)}};

%% ── Nivel 1: Presentación ──────────────────────────────────────────────────
\node[layer=teal] (ui) at (0, 17.8)
  {\textsc{Presentación}\\[3pt]%
   {\normalfont\footnotesize React\;/\;Vite }};

%% ── Nivel 2: Lógica de negocio ─────────────────────────────────────────────
\node[layer=Micolor1] (api) at (0, 15.3)
  {\textsc{Lógica de negocio — FastAPI}\\[3pt]%
   {\normalfont\footnotesize
    Uvicorn\;|\;CORSMiddleware\;|\;
    PyJWT\;+\;bcrypt}};

%% ── Nivel 3: Bloque de Agentes ─────────────────────────────────────────────
\draw[rounded corners=8pt, draw=teal!50, fill=teal!3,
      dashed, line width=1pt]
  (-6.6, 9.5) rectangle (6.6, 13.5);
\node[font=\scriptsize\bfseries\color{teal!70}] at (0, 13.85)
  {Capa de agentes\,/\,LLM\,/\,RAG};

%% Fila superior — agentes sin llamada directa a LLM (teal)
\node[subnode=teal] (ag1) at (-4.4, 12.6)
  {Extracción\\DOCX\,/\,GDoc};
\node[subnode=teal] (ag5) at (0, 12.6)
  {Gestor de\\instrumentos};
\node[subnode=teal] (ag6) at (4.4, 12.6)
  {Sincronización\\GWorkspace};

%% Fila inferior — agentes que usan LLM (naranja)
\node[subnode=orange] (ag2) at (-2.2, 10.6)
  {Evaluación\\inteligente};
\node[subnode=orange] (ag3) at (0, 10.6)
  {Asistencia\\LLM};
\node[subnode=orange] (ag4) at (2.2, 10.6)
  {Búsqueda\\semántica\,(RAG)};

%% ── Fila inferior: Persistencia | OpenAI | Google ─────────────────────────

%% Persistencia — izquierda (azul)
\node[extbox, draw=blue!55, fill=blue!5,
      text width=3.0cm, minimum height=2.2cm] (db) at (-4.5, 7.5)
  {\textsc{Persistencia}\\[4pt]%
   {\normalfont\footnotesize
    SQLAlchemy\,2.0\\SQLite\;·\;hnswlib\,0.8}};

%% OpenAI API — centro (naranja)
\node[extbox, draw=orange!65, fill=orange!5,
      text width=3.2cm, minimum height=2.2cm] (openai) at (0, 7.5)
  {\textsc{OpenAI API}\\[5pt]%
   {\normalfont\footnotesize
    GPT\;·\;completions\\text-embed\;·\;embeddings}};

%% Google Workspace — derecha
\draw[rounded corners=6pt, draw=gray!40, fill=gray!3,
      dashed, line width=0.8pt]
  (2.5, 6.2) rectangle (7.7, 9.0);
\node[font=\scriptsize\bfseries\color{gray!65}] at (5.1, 9.32)
  {Google Workspace};
\node[subcap] (gdrive) at (3.4, 8.2)  {Google\\Drive};
\node[subcap] (gdocs)  at (5.1, 8.2)  {Google\\Docs};
\node[subcap] (gmail)  at (6.8, 8.2)  {Gmail};
\node[font=\scriptsize\color{gray!55}, align=center] at (5.1, 6.8)
  {\texttt{google-api-python-client}};

%% ────────────────── FLECHAS ────────────────────────────────────────────────

%% Cliente ↔ Presentación
\draw[arr] (client.south) --
  node[right, font=\scriptsize\color{gray}, xshift=4pt] {HTTP}
  (ui.north);

%% Presentación ↔ API
\draw[arr] (ui.south) --
  node[right, font=\scriptsize\color{gray}, xshift=4pt] {REST\,/\,JSON}
  (api.north);

%% API ↔ Capa agentes
\draw[arr] (api.south) -- (0, 13.5);

%% API ↔ Persistencia
\draw[arr] (api.south west) to[out=250, in=90] (db.north);

%% Agentes LLM ↔ OpenAI — flechas casi verticales, limpias
\draw[arrdash] (ag2.south) to[out=270, in=125] (openai.north west);
\draw[arrdash] (ag3.south) -- (openai.north);
\draw[arrdash] (ag4.south) to[out=270, in=55] (openai.north east);

%% RAG ↔ Persistencia (hnswlib)
\draw[arrdash] (ag4.south west) to[out=240, in=55] (db.north east);

%% Agentes ↔ Google Workspace
\draw[arrdash] (ag1.south east) to[out=300, in=130] (gdocs.north);
\draw[arrdash] (ag6.south) to[out=270, in=105] (gdrive.north west);

\end{tikzpicture}

\caption{Arquitectura en capas de \textsc{macau}.
  \textit{Fuente: elaboración propia.}}
\label{fig:arquitectura_capas}
\end{figure}

\subsection{Diagrama de despliegue}
\label{subsec:despliegue}

La figura~\ref{fig:despliegue} muestra el diagrama de despliegue de
\textsc{macau}. La arquitectura distingue dos nodos principales: el
\textbf{cliente} (navegador web del usuario) y el \textbf{servidor}, que
aloja simultáneamente el \textit{frontend} y el \textit{backend}.

\begin{figure}[H]
\centering
\scalebox{0.85}{%
\begin{tikzpicture}[
  proc/.style={
    rectangle, rounded corners=5pt,
    draw=#1!70, fill=#1!8, line width=1.2pt,
    font=\small\bfseries, text width=4.6cm,
    align=center, inner sep=10pt, minimum height=2.4cm
  },
  store/.style={
    rectangle, rounded corners=4pt,
    draw=blue!45, fill=blue!5, line width=0.9pt,
    font=\scriptsize, text width=13.0cm,
    align=center, inner sep=8pt
  },
  cloud/.style={
    rectangle, rounded corners=9pt,
    draw=gray!50, fill=gray!5, dashed, line width=0.9pt,
    font=\small\bfseries, text width=4.2cm,
    align=center, inner sep=9pt, minimum height=2.4cm
  },
  arr/.style={
    {Stealth[length=5pt]}-{Stealth[length=5pt]},
    thick, color=gray!65
  },
  extarr/.style={
    {Stealth[length=4pt]}-{Stealth[length=4pt]},
    thick, dashed, color=gray!70
  },
  zonelabel/.style={font=\tiny\bfseries\color{black!55}, anchor=north west,
    fill=white, inner sep=1.5pt}
]

%% ═══════════════════════════════════════════════
%% ZONA 1 — Cliente
%% ═══════════════════════════════════════════════
\draw[rounded corners=8pt, draw=violet!40, fill=violet!5, line width=1.0pt]
  (-3.8, 18.7) rectangle (3.8, 20.8);
\node[zonelabel] at (-3.7, 20.8) {ZONA\,1\;—\;Cliente};
\node[proc=violet, text width=3.2cm, minimum height=1.6cm]
  (browser) at (0, 19.8)
  {Navegador\\[-2pt]{\normalfont\scriptsize (usuario)}};

%% ═══════════════════════════════════════════════
%% ZONA 2 — Máquina local
%% ═══════════════════════════════════════════════
\draw[rounded corners=8pt, draw=black!35, fill=black!2, line width=1.1pt]
  (-7.4, 11.0) rectangle (7.4, 17.8);
\node[zonelabel] at (-7.3, 17.8) {ZONA\,2\;—\;Máquina local (servidor)};

%% Proceso: Frontend
\node[proc=teal] (vite) at (-3.5, 15.5)
  {Frontend\\[4pt]
   {\normalfont\ttfamily\small Puerto\;5173}};

%% Proceso: Backend
\node[proc=Micolor1] (uvicorn) at (3.5, 15.5)
  {Backend\\[4pt]
   {\normalfont\ttfamily\small Puerto\;8011}};

%% Almacenamiento local en disco
\node[store] (storage) at (0, 12.9)
  {\textbf{Disco local}\\[5pt]
   \texttt{macau.db}\qquad
   \texttt{vector\_index/}\qquad
   \texttt{uploads/}\qquad
   \texttt{.env}};

%% ═══════════════════════════════════════════════
%% ZONA 3 — Servicios en la nube (externos)
%% ═══════════════════════════════════════════════
\draw[rounded corners=8pt, draw=orange!45, fill=orange!5, line width=1.0pt]
  (-7.4, 6.5) rectangle (7.4, 10.0);
\node[zonelabel] at (-7.3, 10.0) {ZONA\,3\;—\;Servicios en la nube (externos)};
\node[cloud] (openai) at (-4.0, 8.5)
  {OpenAI API\\[5pt]
   {\normalfont\scriptsize
    GPT-4o\;·\;text-embedding-3\\
    \texttt{api.openai.com}}};

\node[cloud] (google) at (4.0, 8.5)
  {Google Workspace\\[5pt]
   {\normalfont\scriptsize
    Drive\;·\;Docs\;·\;Gmail\\
    OAuth\,2.0}};

%% ── Flechas ────────────────────────────────────────────────────────────────

%% Navegador ↔ Vite (carga de la SPA)
\draw[arr] (browser.south west) to[out=218, in=98]
  node[left, font=\scriptsize\color{gray!80}, xshift=-4pt]
    {\normalfont\texttt{HTTP :5173}}
  (vite.north);

%% Navegador ↔ Uvicorn (llamadas REST)
\draw[arr] (browser.south east) to[out=322, in=82]
  node[right, font=\scriptsize\color{gray!80}, xshift=4pt]
    {\normalfont\texttt{HTTP :8011}}
  (uvicorn.north);

%% Vite ↔ Uvicorn (proxy interno)
\draw[arr] (vite.east) --
  node[above, font=\scriptsize\color{gray!80}, yshift=3pt]
    {proxy\;\texttt{/api}}
  (uvicorn.west);

%% Uvicorn ↔ Disco local
\draw[arr] (uvicorn.south) to[out=260, in=65]
  node[right, font=\scriptsize\color{gray!80}, xshift=3pt]
    {R/W}
  (storage.north);

%% Uvicorn → OpenAI
\draw[extarr] (uvicorn.south)
  to[out=235, in=80]
  node[near end, left, font=\scriptsize\color{gray!70}, xshift=-3pt] {HTTPS/TLS}
  (openai.north);

%% Uvicorn → Google
\draw[extarr] (uvicorn.south)
  to[out=270, in=90]
  node[near end, right, font=\scriptsize\color{gray!70}, xshift=3pt] {HTTPS+OAuth\,2}
  (google.north);

\end{tikzpicture}
}% end scalebox
\caption{Diagrama de despliegue de \textsc{macau}. \textit{Fuente: elaboración propia.}}
\label{fig:despliegue}
\end{figure}

El diagrama organiza el sistema en tres zonas. 

\begin{itemize}
  \item \textbf{Zona\,1: Cliente.} El navegador del usuario, que accede a la
    interfaz web y genera todas las interacciones.
  \item \textbf{Zona\,2: Máquina local.} El servidor que agrupa los procesos que corren en la máquina
local: el \textit{frontend} (puerto 5173) sirve la interfaz;
el \textit{backend} (puerto 8011) expone la API REST, ejecuta los
agentes y gestiona la lógica de negocio; el disco local almacena la
base de datos, el índice vectorial, los archivos subidos y las
variables de entorno. 
  \item \textbf{Zona\,3: Servicios en la nube.} Los servicios externos con los que se comunica el backend: la API de OpenAI (modelos de lenguaje e incrustaciones) y Google Workspace (Drive, Docs y Gmail), ambos accedidos mediante HTTPS.
\end{itemize}


La siguiente estructura de directorios resume la organización del
repositorio:

\begin{quote}
\begin{lstlisting}[basicstyle=\fontsize{9pt}{11pt}\selectfont\ttfamily, frame=none,
    aboveskip=0pt, belowskip=0pt]
TesisMCD/
  backend/
    app/
      main.py          <- enrutador principal (FastAPI)
      models.py        <- entidades SQLAlchemy (ORM)
      schemas.py       <- esquemas de validacion (Pydantic)
      auth.py          <- autenticacion JWT + bcrypt
      database.py      <- sesion y motor SQLAlchemy
      docx_import.py   <- importacion de propuestas DOCX
      docx_export.py   <- exportacion de propuestas DOCX
    agents/
      extract.py       <- extraccion de propuestas DOCX/GDoc
      evaluation.py    <- evaluacion inteligente con LLM
      assistance.py    <- asistencia y reformulacion LLM
      search.py        <- busqueda semantica RAG (hnswlib)
      instruments.py   <- gestion de instrumentos evaluativos
      gsync.py         <- sincronizacion con Google Workspace
    data/
      macau.db         <- base de datos SQLite principal
      uploads/         <- archivos subidos (propuestas, evidencias)
      vector_index/    <- indice hnswlib (embeddings semanticos)
  frontend/
    src/
      App.jsx          <- componente raiz y enrutamiento
      main.jsx         <- punto de entrada React
      styles.css       <- estilos globales
    index.html         <- shell HTML
  docs/
    ARCHITECTURE.md    <- descripcion de la arquitectura
    DECISIONS.md       <- registro de decisiones tecnicas
  .env                 <- claves de API y variables de entorno
\end{lstlisting}
\end{quote}

%% =============================================================================
\section{Modelo de datos}
\label{sec:modelo_datos}
%% =============================================================================

El modelo de datos de \textsc{macau} se implementa con SQLAlchemy~2.0 sobre
SQLite, declarado en \texttt{app/models.py} mediante el patrón ORM declarativo.
La elección de SQLite responde a los requisitos concretos de despliegue: el
sistema opera en una única máquina institucional sin infraestructura de base de
datos dedicada, con concurrencia baja y volumen de datos acotado. SQLAlchemy
abstrae el motor subyacente, de modo que una eventual migración a PostgreSQL
requeriría únicamente modificar la cadena de conexión en \texttt{database.py},
sin cambios en el código de negocio.

El esquema se organiza en seis grupos funcionales con responsabilidades claramente
delimitadas, que se ilustran en la figura~\ref{fig:er_grupos}.

%% Layout circular: Propuestas al centro, 5 grupos en corona a radio=5.4cm
%% Ángulos (sentido horario desde arriba):
%%   Plan de estudios   →  72°  → ( 5.13,  1.67)
%%   Controles          →  144° → (-3.17,  4.37)  (NO: usamos sentido antihorario)
%% Convención antihoraria desde las 12h:
%%   doc   →  -126° = 234°  → (-4.37, -3.17)  NO
%% Usamos ángulos en grados desde eje +x, antihorario:
%%   plan:  18°   → (5.13,  1.67)   derecha-arriba
%%   instr: -54°  → (3.17, -4.37)   derecha-abajo
%%   acc:  -126°  → (-3.17,-4.37)   izquierda-abajo
%%   doc:   162°  → (-5.13, 1.67)   izquierda-arriba
%%   ctrl:   90°  → (0,     5.40)   arriba
\begin{figure}[H]
\centering
\scalebox{0.82}{%
\begin{tikzpicture}[
  grp/.style={
    rectangle, rounded corners=6pt,
    line width=1.1pt,
    font=\small\bfseries, text width=4.4cm,
    align=center, inner sep=10pt, minimum height=1.9cm
  },
  brel/.style={
    {Stealth[length=4.5pt]}-{Stealth[length=4.5pt]},
    thick, color=gray!50
  },
  rel/.style={
    -{Stealth[length=4.5pt]},
    thick, color=gray!50
  },
  rlabel/.style={
    font=\scriptsize\color{gray!65},
    fill=white, inner sep=2pt, rounded corners=2pt
  }
]

%% ── Radio de la corona ───────────────────────────────────────────────────
%% r = 5.4;  pos en (r·cos θ, r·sin θ)
%% ctrl:  90°  → (0,      5.40)
%% plan:  18°  → (5.130,  1.669)
%% instr:-54°  → (3.173, -4.369)
%% acc: -126°  → (-3.173,-4.369)
%% doc:  162°  → (-5.130, 1.669)

%% ── Nodo central ─────────────────────────────────────────────────────────
\node[grp, draw=Micolor1!70, fill=Micolor1!8] (prop) at (0, 0)
  {\textsc{Propuestas}\\[4pt]
   {\normalfont\scriptsize
    \texttt{Proposal}\\
    \texttt{ProposalTeacher}\\
    \texttt{ProposalCompetency}\\
    \texttt{Proposal...ControlResult}}};

%% ── Nodos periféricos ────────────────────────────────────────────────────
\node[grp, draw=orange!70, fill=orange!8] (ctrl) at (0, 5.40)
  {\textsc{Controles inteligentes}\\[4pt]
   {\normalfont\scriptsize
    \texttt{IntelligentControl}\\
    \texttt{IntelligentCtrlSettings}}};

\node[grp, draw=blue!70, fill=blue!8] (plan) at (6.2, 0)
  {\textsc{Plan de estudios}\\[4pt]
   {\normalfont\scriptsize
    \texttt{StudyPlan/Year/Term}\\
    \texttt{StudySubject}\\
    \texttt{StudySubjectPrereq.}\\
    \texttt{CompetencyCatalog}}};

\node[grp, draw=violet!70, fill=violet!8] (instr) at (3.2, -4.6)
  {\textsc{Instrumentos evaluativos}\\[4pt]
   {\normalfont\scriptsize
    \texttt{EvaluativeInstrument}\\
    \texttt{EvaluativeInstrumentFolder}}};

\node[grp, draw=green!60!black, fill=green!10] (acc) at (-3.2, -4.6)
  {\textsc{Acreditación}\\[4pt]
   {\normalfont\scriptsize
    \texttt{AccreditationSettings}\\
    \texttt{EvidenceRegistry}\\
    \texttt{EvidenceVersion/AuditLog}\\
    \texttt{WorkPlanActivity/Task}}};

\node[grp, draw=teal!70, fill=teal!8] (doc) at (-6.2, 0)
  {\textsc{Docentes}\\[4pt]
   {\normalfont\scriptsize
    \texttt{Teacher}\\
    \texttt{TeacherCareer}\\
    \texttt{Career}\\
    \texttt{CareerCommitteeMember}}};

%% ── Relaciones radiales (prop ↔ periferia) ───────────────────────────────
%% Cada arista va del centro al nodo periférico: sin cruces garantizados.

\draw[rel] (prop) -- node[rlabel, pos=0.52, sloped, above=1pt] {resultados LLM} (ctrl);
\draw[rel] (prop) -- node[rlabel, pos=0.52, sloped, above=1pt] {FK asignatura}  (plan);
\draw[brel](prop) -- node[rlabel, pos=0.52, sloped, above=1pt] {evidencia RA}    (instr);
\draw[brel](prop) -- node[rlabel, pos=0.52, sloped, above=1pt] {marco}           (acc);
\draw[brel](prop) -- node[rlabel, pos=0.52, sloped, above=1pt] {N:M docentes}    (doc);

%% ── Relaciones periféricas ────────────────────────────────────────────────
%% plan → instr: arco corto entre nodos adyacentes (72° de separación)
\draw[rel]  (plan.south) to[out=270, in=30]
  node[rlabel, pos=0.45, right=2pt] {asignatura} (instr.north east);

%% doc ↔ instr: arco largo entre nodos opuestos (216° de separación)
\draw[brel] (doc.south) to[out=270, in=210]
  node[rlabel, pos=0.5, below=3pt] {autoría} (instr.west);

%% doc ↔ acc: arco corto entre nodos adyacentes
\draw[brel] (doc.south east) to[out=310, in=150]
  node[rlabel, pos=0.5, below=2pt] {participación} (acc.north west);

%% plan ↔ acc: arco largo por la parte exterior inferior
\draw[brel] (plan.south) to[out=250, in=30]
  node[rlabel, pos=0.5, right=2pt] {marco} (acc.north east);

\end{tikzpicture}}% end scalebox
\caption{Grupos funcionales del modelo de datos. \textit{Fuente: elaboración propia.}}
\label{fig:er_grupos}
\end{figure}

\subsubsection*{Grupo 1 — Núcleo de propuestas}

La entidad central del sistema es \texttt{Proposal}. Su diseño consolida en una
única tabla todos los campos que conforman una propuesta de asignatura, organizados
en cuatro categorías funcionales: \textit{identificación académica},
\textit{carga horaria}, \textit{contenido pedagógico} y
\textit{control e integración externa}. El cuadro~\ref{tab:proposal} detalla los
atributos principales de la entidad.

\begin{longtable}{@{}p{4.2cm} p{2.0cm} p{7.6cm}@{}}
\caption{Atributos principales de la entidad \texttt{Proposal}.}
\label{tab:proposal}\\
\toprule
\textbf{Atributo} & \textbf{Tipo} & \textbf{Descripción} \\
\midrule
\endfirsthead
\multicolumn{3}{@{}l}{\small\textit{(continuación de la tabla~\ref{tab:proposal})}}\\[4pt]
\toprule
\textbf{Atributo} & \textbf{Tipo} & \textbf{Descripción} \\
\midrule
\endhead
\midrule
\multicolumn{3}{r@{}}{\small\textit{continúa en la página siguiente}}\\
\endfoot
\bottomrule
\endlastfoot
\multicolumn{3}{@{}l}{\textit{Identificación académica}} \\[2pt]
\texttt{title}            & String  & Denominación oficial de la asignatura. \\
\texttt{career}           & String  & Carrera a la que pertenece la propuesta. \\
\texttt{subject}          & String  & Nombre normalizado de la asignatura. \\
\texttt{study\_plan}      & String  & Plan de estudios de referencia. \\
\texttt{quarter}          & String  & Cuatrimestre o período lectivo. \\
\texttt{character}        & String  & Carácter: obligatoria u optativa. \\
\texttt{regime}           & String  & Régimen: cuatrimestral o anual. \\
\midrule
\multicolumn{3}{@{}l}{\textit{Carga horaria}} \\[2pt]
\texttt{theoretical\_hours} & Integer & Horas teóricas semanales. \\
\texttt{practical\_hours}   & Integer & Horas prácticas semanales. \\
\texttt{total\_hours}       & Integer & Total de horas de la asignatura. \\
\midrule
\multicolumn{3}{@{}l}{\textit{Contenido pedagógico}} \\[2pt]
\texttt{minimum\_content}   & Text    & Contenidos mínimos según plan. \\
\texttt{fundamentals\_part1}& Text    & Fundamentación: importancia de la asignatura. \\
\texttt{fundamentals\_part2}& Text    & Fundamentación: perfil profesional. \\
\texttt{methodology}        & Text    & Metodología de enseñanza. \\
\texttt{evaluation}         & Text    & Sistema de evaluación. \\
\texttt{bibliography}       & Text    & Bibliografía básica y complementaria. \\
\texttt{learning\_outcomes} & JSON    & Lista estructurada de resultados de aprendizaje. \\
\texttt{units}              & JSON    & Unidades temáticas con horas por unidad. \\
\texttt{practicals}         & JSON    & Trabajos prácticos de la asignatura. \\
\texttt{teaching\_team}     & JSON    & Equipo docente con roles y titulación. \\
\midrule
\multicolumn{3}{@{}l}{\textit{Control e integración externa}} \\[2pt]
\texttt{status}             & String  & Estado del ciclo de vida de la propuesta. \\
\texttt{editing\_locked}    & Boolean & Bloqueo de edición activado por el directivo. \\
\texttt{gdoc\_url}          & String  & URL del Google Doc vinculado. \\
\texttt{gdoc\_last\_synced} & DateTime& Marca temporal de la última sincronización. \\
\texttt{gdoc\_status}       & String  & Estado de la sincronización (\textit{synced}, \textit{outdated}, etc.). \\
\texttt{source\_type}       & String  & Origen: \texttt{docx}, \texttt{gdoc} o \texttt{manual}. \\
\texttt{created\_at}        & DateTime& Fecha y hora de creación del registro. \\
\texttt{updated\_at}        & DateTime& Fecha y hora de la última modificación. \\
\end{longtable}

Las secciones de contenido de estructura variable ---resultados de aprendizaje,
unidades temáticas, trabajos prácticos y equipo docente--- se almacenan como
columnas de tipo \texttt{JSON}. Esta decisión permite representar colecciones de
longitud arbitraria sin normalizar en tablas adicionales, a costa de renunciar a
consultas relacionales sobre esos campos; dado que el acceso a dichas secciones
siempre se realiza de forma completa (lectura y escritura de la propuesta entera),
el enfoque es apropiado en este contexto.

\textbf{\textit{Entidades satélite del grupo:}} \texttt{ProposalTeacher} materializa la relación N:M entre propuestas y docentes,
permitiendo que múltiples docentes participen en una misma propuesta con roles
diferenciados. \texttt{ProposalCompetency} vincula cada propuesta con las
competencias CONEAU que declara cubrir, tomadas del catálogo institucional.
\texttt{ProposalIntelligentControlResult} persiste el diagnóstico producido por
cada control normativo automático: registra si la verificación fue superada,
qué aspecto falló, por qué razón y qué texto alternativo propuso el modelo de
lenguaje; estos datos alimentan la vista de retroalimentación del docente y el
panel de auditoría del directivo.

La especificación detallada de este grupo figura en \textbf{\textit{Anexo~\ref{anexo:modelo_datos} — Especificación del modelo de datos, Grupo 1}} (véase pág.~\pageref{anexo:grupo1_modelo_datos}).

\subsubsection*{Grupo 2 — Gestión docente}

La entidad central de este grupo es \texttt{Teacher}, que consolida en una
sola tabla la identidad institucional del docente y sus credenciales de
autenticación. Este diseño evita tablas \texttt{User} separadas: cada docente
es simultáneamente un actor del sistema con sus propios permisos. La
tabla~\ref{tab:teacher} recoge sus atributos completos.

\begin{longtable}{@{}p{4.0cm} p{2.0cm} p{7.6cm}@{}}
\caption{Entidad \texttt{teachers} — perfil e identidad docente.}
\label{tab:teacher}\\
\toprule
\textbf{Atributo} & \textbf{Tipo} & \textbf{Descripción} \\
\midrule
\endfirsthead
\multicolumn{3}{@{}l}{\small\textit{(continuación — \texttt{teachers})}}\\[3pt]
\toprule
\textbf{Atributo} & \textbf{Tipo} & \textbf{Descripción} \\
\midrule
\endhead
\midrule\multicolumn{3}{r@{}}{\small\textit{continúa en la página siguiente}}\\
\endfoot
\bottomrule
\endlastfoot
%% ── Identificación y perfil ────────────────────────
\multicolumn{3}{@{}l}{\textit{Identificación y perfil}} \\
\midrule
\texttt{id}             & Integer  & Identificador interno autoincremental (PK). \\
\texttt{name}           & String   & Nombre completo del docente. \\
\texttt{normalized\_key}& String   & Clave normalizada para búsqueda sin diferencias
                                      tipográficas (tildes, mayúsculas, espacios). \\
\texttt{email}          & String   & Correo electrónico institucional. \\
\texttt{category}       & String   & Categoría docente (p.\,ej.\ titular, adjunto, JTP). \\
\texttt{dedication}     & String   & Régimen de dedicación (exclusiva, semi, simple). \\
\midrule
%% ── Autenticación y acceso ─────────────────────────
\multicolumn{3}{@{}l}{\textit{Autenticación y acceso}} \\
\midrule
\texttt{password\_hash}        & String   & Hash bcrypt de la contraseña local.
                                            Nulo si el docente no activó
                                            autenticación por credencial. \\
\texttt{last\_login}           & DateTime & Marca temporal del último acceso al sistema. \\
\texttt{password\_reset\_at}   & DateTime & Marca temporal de la última solicitud de
                                            restablecimiento de contraseña. \\
\texttt{gmail\_refresh\_token} & Text     & Token de refresco OAuth~2.0 para la integración
                                            con Google Drive del docente. \\
\texttt{is\_admin}            & Boolean  & Indica si el docente posee privilegios de
                                            administración global del sistema. \\
\midrule
%% ── Control temporal ───────────────────────────────
\multicolumn{3}{@{}l}{\textit{Control temporal}} \\
\midrule
\texttt{created\_at} & DateTime & Fecha y hora de creación (valor por defecto: ahora). \\
\texttt{updated\_at} & DateTime & Fecha y hora de última modificación (auto actualizado). \\
\end{longtable}

\textit{\textbf{Entidades satélite del grupo:}} \texttt{TeacherCareer} registra la pertenencia
de un docente a una carrera mediante un par \texttt{teacher\_id}
$\times$~\texttt{career} (nombre de carrera como cadena). \texttt{Career}
centraliza el catálogo de carreras activas con referencias opcionales a los
docentes que ejercen como director (\texttt{director\_id}) y secretario
(\texttt{secretario\_id}). \texttt{CareerCommitteeMember} materializa la
comisión curricular como una relación N:M entre \texttt{Career} y
\texttt{Teacher} con eliminación en cascada cuando se elimina la carrera o
el docente.

La especificación detallada de este grupo figura en \textbf{\textit{Anexo~\ref{anexo:modelo_datos} — Especificación del modelo de datos, Grupo 2}} (véase pág.~\pageref{anexo:grupo2_modelo_datos}).

\subsubsection*{Grupo 3 — Plan de estudios}

El modelo curricular se organiza como una jerarquía de cuatro niveles:
\texttt{StudyPlan} $\to$ \texttt{StudyYear} $\to$ \texttt{StudyTerm} $\to$
\texttt{StudySubject}. Esta estructura replica fielmente el esquema formal de
la carrera y permite que cada asignatura quede unívocamente ubicada en su año
y cuatrimestre dentro de un plan vigente. La entidad hoja de la jerarquía,
\texttt{StudySubject}, actúa como ancla para las propuestas pedagógicas
mediante la clave foránea \texttt{study\_subject\_id} de \texttt{Proposal};
sus atributos se detallan en la tabla~\ref{tab:study_subject}.

\begin{longtable}{@{}p{4.0cm} p{2.0cm} p{7.6cm}@{}}
\caption{Entidad \texttt{study\_subjects} — asignatura del plan de estudios.}
\label{tab:study_subject}\\
\toprule
\textbf{Atributo} & \textbf{Tipo} & \textbf{Descripción} \\
\midrule
\endfirsthead
\multicolumn{3}{@{}l}{\small\textit{(continuación — \texttt{study\_subjects})}}\\[3pt]
\toprule
\textbf{Atributo} & \textbf{Tipo} & \textbf{Descripción} \\
\midrule
\endhead
\midrule\multicolumn{3}{r@{}}{\small\textit{continúa en la página siguiente}}\\
\endfoot
\bottomrule
\endlastfoot
\multicolumn{3}{@{}l}{\textit{Identificación}} \\
\midrule
\texttt{id}                   & Integer & Identificador interno autoincremental (PK). \\
\texttt{term\_id}             & Integer & FK a \texttt{study\_terms.id}; ubica la asignatura en su cuatrimestre. \\
\texttt{code}                 & String  & Código oficial de la asignatura en el plan. \\
\texttt{name}                 & String  & Denominación oficial de la asignatura. \\
\midrule
\multicolumn{3}{@{}l}{\textit{Carga horaria y régimen}} \\
\midrule
\texttt{character}            & String  & Carácter: \textit{Obligatoria} u \textit{Optativa}. \\
\texttt{regime}               & String  & Régimen: \textit{Cuatrimestral} o \textit{Anual}. \\
\texttt{theoretical\_hours}  & Integer & Horas teóricas semanales. \\
\texttt{practical\_hours}    & Integer & Horas prácticas semanales. \\
\texttt{total\_hours}        & Integer & Total de horas de la asignatura. \\
\texttt{weekly\_hours}       & Integer & Carga horaria semanal combinada. \\
\texttt{practice\_scope}     & Text    & Alcance del trabajo práctico según el plan. \\
\midrule
\multicolumn{3}{@{}l}{\textit{Contenido curricular}} \\
\midrule
\texttt{minimum\_content}       & Text & Contenidos mínimos establecidos por el plan. \\
\texttt{generic\_competencies}  & Text & Competencias genéricas CONEAU referenciadas. \\
\texttt{specific\_competencies} & Text & Competencias específicas CONEAU referenciadas. \\
\texttt{blocks}                 & JSON & Bloques temáticos con horas asignadas por bloque. \\
\midrule
\multicolumn{3}{@{}l}{\textit{Control temporal}} \\
\midrule
\texttt{created\_at} & DateTime & Fecha y hora de creación (\textit{server default}: ahora). \\
\texttt{updated\_at} & DateTime & Fecha y hora de última modificación (auto actualizado). \\
\end{longtable}

\textbf{\textit{Entidades complementarias del grupo:}} \texttt{StudyPlan} encabeza la jerarquía
con el nombre del plan y un campo \texttt{payload} (JSON) para metadatos
extensos sin esquema fijo. \texttt{StudyYear} y \texttt{StudyTerm} son nodos
intermedios de ordenamiento con campos \texttt{sort\_order} que permiten
respetar el orden oficial sin depender del valor numérico del \texttt{id}.
\texttt{StudySubjectPrerequisite} expresa las correlatividades como una
autorreferencia sobre \texttt{study\_subjects}: cada fila indica que
\texttt{subject\_id} requiere aprobar \texttt{prerequisite\_id}.
\texttt{CompetencyCatalog} centraliza el catálogo de competencias CONEAU por
carrera y plan, clasificadas como genéricas o específicas mediante
\texttt{competency\_type}; sirve como fuente de verdad para poblar las
competencias de nuevas propuestas.

La especificación detallada de este grupo figura en \textbf{\textit{Anexo~\ref{anexo:modelo_datos} — Especificación del modelo de datos, Grupo 3}} (véase pág.~\pageref{anexo:grupo3_modelo_datos}).

\subsubsection*{Grupo 4 — Controles inteligentes}

Este grupo implementa el motor de verificación normativa automática.
\texttt{IntelligentControl} define cada regla verificable: un \texttt{topic}
categorizador, un \texttt{name} descriptivo y el campo \texttt{instruction},
que contiene el prompt completo que se envía al LLM junto con el texto de la
propuesta. El campo \texttt{associated\_topics} (JSON) permite agrupar
controles relacionados para ejecuciones consolidadas. La
tabla~\ref{tab:intelligent_control} recoge sus atributos.

\begin{longtable}{@{}p{4.0cm} p{2.0cm} p{7.6cm}@{}}
\caption{Entidad \texttt{intelligent\_controls} — definición de controles LLM.}
\label{tab:intelligent_control}\\
\toprule
\textbf{Atributo} & \textbf{Tipo} & \textbf{Descripción} \\
\midrule
\endfirsthead
\multicolumn{3}{@{}l}{\small\textit{(continuación — \texttt{intelligent\_controls})}}\\[3pt]
\toprule
\textbf{Atributo} & \textbf{Tipo} & \textbf{Descripción} \\
\midrule
\endhead
\bottomrule
\endlastfoot
\texttt{id}                 & Integer & Identificador interno autoincremental (PK). \\
\texttt{topic}              & String  & Categoría temática del control (p.\,ej.\ \textit{competencias}, \textit{bibliografía}). \\
\texttt{name}               & String  & Nombre descriptivo del control. \\
\texttt{instruction}        & Text    & Prompt completo enviado al LLM junto con el texto de la propuesta. \\
\texttt{is\_active}        & Boolean & Indica si el control participa en las ejecuciones (defecto: verdadero). \\
\texttt{sort\_order}       & Integer & Orden de presentación en la interfaz. \\
\texttt{associated\_topics}& JSON    & Lista de tópicos relacionados para ejecuciones consolidadas. \\
\texttt{created\_at}       & DateTime& Fecha y hora de creación (\textit{server default}: ahora). \\
\texttt{updated\_at}       & DateTime& Fecha y hora de última modificación (auto actualizado). \\
\end{longtable}

\medskip
\noindent\textbf{Configuración de modos de ejecución.}
La entidad \texttt{IntelligentControlSettings} centraliza los par\'{a}metros
de los tres perfiles disponibles. Cada modo tiene su propio modelo de
lenguaje, temperatura y l\'{i}mite de \textit{tokens}, todos configurables
de forma independiente sin tocar el c\'{o}digo:

\begin{description}[noitemsep, leftmargin=1.4cm, labelwidth=1.2cm]
  \item[\textit{Guepardo}] R\'{a}pido y de bajo costo.
    Modelo por defecto: \texttt{gpt-4o-mini}, temperatura~0{,}15,
    m\'{a}x.\ 420~\textit{tokens}.
    Perfil asignado por defecto a \textbf{docentes}
    (\texttt{docente\_mode}).
  \item[\textit{Delf\'{i}n}] Equilibrado entre velocidad y profundidad.
    Modelo por defecto: \texttt{gpt-4o-mini}, temperatura~0{,}10,
    m\'{a}x.\ 500~\textit{tokens}.
    Perfil asignado por defecto a \textbf{directivos}
    (\texttt{director\_last\_mode}).
  \item[\textit{Ballena}] M\'{a}xima profundidad de an\'{a}lisis.
    Modelo por defecto: \texttt{gpt-4o}, temperatura~0{,}10,
    m\'{a}x.\ 700~\textit{tokens}.
    
\end{description}

\noindent La selecci\'{o}n activa se persiste por rol:
\texttt{director\_last\_mode} para directivos y
\texttt{docente\_mode} para docentes, de modo que cada usuario retoma
el \'{u}ltimo perfil utilizado en su pr\'{o}xima sesi\'{o}n.
Un perfil no v\'{a}lido se normaliza autom\'{a}ticamente a
\textit{Delf\'{i}n}.

Los valores indicados son los \textbf{par\'{a}metros por defecto}.
Un perfil con rol directivo puede modificarlos en cualquier momento a
trav\'{e}s del endpoint \texttt{PATCH /intelligent-controls/settings}:
modelo de lenguaje, temperatura y l\'{i}mite de \textit{tokens} se
pueden ajustar de forma independiente para cada uno de los tres modos,
sin necesidad de intervenci\'{o}n técnica en el servidor.
Esto permite, por ejemplo, migrar \textit{Ballena} a un modelo más
reciente cuando est\'{e} disponible, o subir el l\'{i}mite de
\textit{tokens} de \textit{Guepardo} durante períodos de mayor
exigencia de an\'{a}lisis, sin tocar el c\'{o}digo fuente.

La especificación detallada de este grupo figura en \textbf{\textit{Anexo~\ref{anexo:modelo_datos} — Especificación del modelo de datos, Grupo 4}} (véase pág.~\pageref{anexo:grupo4_modelo_datos}).

\subsubsection*{Grupo 5 — Instrumentos evaluativos}

La gestión de evaluaciones se organiza en dos entidades. \texttt{EvaluativeInstrument}
registra trabajos prácticos, parciales, finales y recuperatorios de una
asignatura, discriminados por \texttt{instrument\_type}. Soporta dos modalidades
de almacenamiento: (a) local (\texttt{stored\_filename}, \texttt{file\_path})
o (b) en Google Drive (\texttt{gdrive\_url}, \texttt{gdrive\_file\_id}), lo que
permite tanto la carga directa de archivos como la vinculación con documentos
ya existentes en el espacio Drive institucional. La tabla~\ref{tab:evaluative_instrument}
detalla sus atributos.

\begin{longtable}{@{}p{4.0cm} p{2.0cm} p{7.6cm}@{}}
\caption{Entidad \texttt{evaluative\_instruments} — registro de evaluaciones.}
\label{tab:evaluative_instrument}\\
\toprule
\textbf{Atributo} & \textbf{Tipo} & \textbf{Descripción} \\
\midrule
\endfirsthead
\multicolumn{3}{@{}l}{\small\textit{(continuación — \texttt{evaluative\_instruments})}}\\[3pt]
\toprule
\textbf{Atributo} & \textbf{Tipo} & \textbf{Descripción} \\
\midrule
\endhead
\bottomrule
\endlastfoot
\multicolumn{3}{@{}l}{\textit{Identificación y contexto}} \\
\midrule
\texttt{id}                & Integer & Identificador interno autoincremental (PK). \\
\texttt{career}            & String  & Carrera a la que pertenece la evaluación. \\
\texttt{study\_plan}       & String  & Plan de estudios de referencia. \\
\texttt{subject}           & String  & Asignatura a la que pertenece la evaluación. \\
\texttt{instrument\_type}  & String  & Tipo: \textit{TP} \textbar{} \textit{Parcial} \textbar{} \textit{Final} \textbar{} \textit{Recuperatorio}. \\
\texttt{title}             & String  & Denominación descriptiva del instrumento (opcional). \\
\midrule
\multicolumn{3}{@{}l}{\textit{Almacenamiento local}} \\
\midrule
\texttt{original\_filename} & String  & Nombre original del archivo subido por el usuario. \\
\texttt{stored\_filename}   & String  & Nombre del archivo almacenado en el servidor. \\
\texttt{file\_path}        & String  & Ruta del archivo en el servidor (ruta relativa o absoluta). \\
\texttt{file\_size}        & Integer & Tamaño del archivo en bytes. \\
\texttt{mime\_type}        & String  & Tipo MIME del archivo (p.\,ej.\ \texttt{application/pdf}). \\
\midrule
\multicolumn{3}{@{}l}{\textit{Almacenamiento en Drive}} \\
\midrule
\texttt{gdrive\_url}       & String  & URL pública o compartida del archivo en Google Drive. \\
\texttt{gdrive\_file\_id}  & String  & Identificador único del archivo en Google Drive API. \\
\midrule
\multicolumn{3}{@{}l}{\textit{Metadata y auditoría}} \\
\midrule
\texttt{uploaded\_by}      & String  & Usuario o email que cargó el documento. \\
\texttt{created\_at}       & DateTime& Fecha y hora de subida (\textit{server default}: ahora). \\
\texttt{updated\_at}       & DateTime& Fecha y hora de última modificación (auto actualizado). \\
\end{longtable}

\textbf{\textit{Carpetas de Drive centralizado:}} \texttt{EvaluativeInstrumentFolder}
asocia una carpeta de Google Drive a una determinada combinación
carrera~+~plan~+~asignatura. Esta estructura permite consolidar todos los
instrumentos de una asignatura bajo un único directorio en Drive, facilitando
la descarga y sincronización programada.

La especificación detallada de este grupo figura en \textbf{\textit{Anexo~\ref{anexo:modelo_datos} — Especificación del modelo de datos, Grupo 5}} (véase pág.~\pageref{anexo:grupo5_modelo_datos}).

\subsubsection*{Grupo 6 — Acreditación}

Este grupo concentra seis entidades de soporte al proceso CONEAU.
\texttt{AccreditationSettings} actúa como entidad cabecera de configuración por
carrera y plan: define carpetas de origen y destino en Drive, modo de
procesamiento, alcance del escaneo y listas JSON de tipos de evidencia y
actores habilitados. La tabla~\ref{tab:accreditation_settings} resume su
estructura.

\begin{longtable}{@{}p{4.0cm} p{2.0cm} p{7.6cm}@{}}
\caption{Entidad \texttt{accreditation\_settings} — parámetros del proceso de acreditación.}
\label{tab:accreditation_settings}\\
\toprule
\textbf{Atributo} & \textbf{Tipo} & \textbf{Descripción} \\
\midrule
\endfirsthead
\multicolumn{3}{@{}l}{\small\textit{(continuación — \texttt{accreditation\_settings})}}\\[3pt]
\toprule
\textbf{Atributo} & \textbf{Tipo} & \textbf{Descripción} \\
\midrule
\endhead
\bottomrule
\endlastfoot
\multicolumn{3}{@{}l}{\textit{Identificación y alcance}} \\
\midrule
\texttt{id}                  & Integer  & Identificador interno autoincremental (PK). \\
\texttt{career}              & String   & Carrera sobre la que aplica la configuración (índice). \\
\texttt{study\_plan}         & String   & Plan de estudios de referencia (opcional, indexado). \\
\midrule
\multicolumn{3}{@{}l}{\textit{Conectividad y procesamiento}} \\
\midrule
\texttt{source\_folder\_url}      & String   & URL de carpeta de origen para relevamiento de evidencias. \\
\texttt{destination\_folder\_url} & String   & URL de carpeta de destino para consolidación documental. \\
\texttt{process\_mode}            & String   & Estrategia de procesamiento: \textit{move} o \textit{copy}. \\
\texttt{recursive\_scan}          & Boolean  & Indica si el relevamiento recorre subcarpetas de forma recursiva. \\
\midrule
\multicolumn{3}{@{}l}{\textit{Parámetros dinámicos}} \\
\midrule
\texttt{evidence\_types} & JSON & Tipos de evidencia habilitados para el proceso acreditador. \\
\texttt{actor\_roles}    & JSON & Roles institucionales admitidos en el flujo de trabajo. \\
\texttt{actors}          & JSON & Actores concretos participantes (equipo operativo). \\
\midrule
\multicolumn{3}{@{}l}{\textit{Control temporal}} \\
\midrule
\texttt{created\_at} & DateTime & Fecha y hora de creación (\textit{server default}: ahora). \\
\texttt{updated\_at} & DateTime & Fecha y hora de última modificación (auto actualizado). \\
\end{longtable}

\textbf{\textit{Entidades complementarias del grupo:}} \texttt{AccreditationEvidenceRegistry}
es el registro principal de evidencias (origen local, URL de Drive o carpeta de
Drive), con trazabilidad de integridad mediante \texttt{checksum\_sha256} y
metadatos JSON. \texttt{AccreditationEvidenceVersion} preserva el historial de
versionado por evidencia, mientras \texttt{AccreditationEvidenceAuditLog}
registra acciones inmutables de creación, actualización y versionado junto con
actor y campos modificados. \texttt{AccreditationWorkPlanActivity} y
\texttt{AccreditationWorkPlanTask} modelan el plan operativo del proceso:
etapas, sub-etapas, responsables, fechas de compromiso, estado y observaciones,
con desagregación en tareas de seguimiento.

La especificación detallada de este grupo figura en \textbf{\textit{Anexo~\ref{anexo:modelo_datos} — Especificación del modelo de datos, Grupo 6}} (véase pág.~\pageref{anexo:grupo6_modelo_datos}).

%% =============================================================================
\section{Módulo de gestión del ciclo de vida de propuestas}
\label{sec:propuestas}
%% =============================================================================

La propuesta de asignatura es la unidad documental central de \textsc{macau}:
agrupa en un solo objeto persistente toda la información exigida —desde los datos de identificación de la asignatura hasta la
bibliografía, la metodología y los resultados de aprendizaje— y la expone a
flujos de validación normativa, revisión colaborativa y exportación
institucional.

El ciclo de vida de una propuesta está gobernado por cinco estados formales
con transiciones explícitas y controladas por rol:
\begin{description}[noitemsep, topsep=3pt]
    \item[\textit{Borrador}:] la propuesta fue creada o importada por el directivo; el docente asignado puede editarla libremente y agregar toda la información requerida, ya que todos sus campos son modificables.
    \item[\textit{En revisión}:] fue enviada a la Comisión Curricular; puede
          quedar bloqueada para edición según la configuración del proceso.
    \item[\textit{Observada}:] la Comisión Curricular detectó no conformidades;
          el docente recibe las observaciones y debe subsanarlas antes de
          volver a enviarla.
    \item[\textit{Aprobada}:] el director aprobó formalmente la propuesta;
          la edición queda bloqueada para preservar la integridad del documento
          de referencia.
    \item[\textit{Archivada}:] cierre de ciclo; la propuesta queda como
          evidencia documental accesible al proceso de acreditación CONEAU.
\end{description}

La Figura~\ref{fig:flujo_estados_propuesta} resume las transiciones válidas
entre estados, incluyendo el ciclo de observación y corrección previo a la
aprobación final.

\begin{figure}[H]
\centering
\shorthandoff{<>}
\begin{tikzpicture}[
  every node/.style={font=\small},
  estado/.style={
    rectangle, rounded corners=4pt,
    draw=Micolor1, line width=0.8pt,
    minimum width=3.4cm, minimum height=0.9cm,
    align=center, fill=Micolor1!12
  },
  decision/.style={
    diamond, aspect=2.4,
    draw=Micolor1, line width=0.8pt,
    minimum height=1.0cm,
    align=center, fill=Micolor1!6,
    font=\small
  },
  actor/.style={font=\scriptsize\itshape, text=gray!70!black},
  etiq/.style={font=\scriptsize\bfseries, fill=white, inner sep=1pt}
]
  % --- Nodos del eje principal ---
  \node[estado]                           (borrador)  {Borrador};
  \node[estado,  below=1.0cm of borrador] (revision)  {En revisi\'{o}n};
  \node[decision, below=1.3cm of revision] (dec)      {¿Conforme?};
  \node[estado,  below=1.4cm of dec]      (aprobada)  {Aprobada};
  \node[estado,  below=1.0cm of aprobada] (archivada) {Archivada};

  % --- Nodo lateral Observada, alineado con el rombo ---
  \node[estado, right=3.6cm of dec]       (observada) {Observada};

  % --- Eje principal ---
  \draw[->, thick, color=Micolor1]
    (borrador) -- node[etiq, right=2pt]{env\'ia}
                  node[actor, left=2pt]{\textit{(Director)}} (revision);

  \draw[->, thick, color=Micolor1]
    (revision) -- node[etiq, right=2pt]{revisa}
                  node[actor, left=2pt]{\textit{(Com.\,Curricular)}} (dec);

  \draw[->, thick, color=Micolor1]
    (dec) -- node[etiq, right=2pt]{S\'i --- aprueba}
              node[actor, left=2pt]{\textit{(Director)}} (aprobada);

  \draw[->, thick, color=Micolor1]
    (aprobada) -- node[etiq, right=2pt]{archiva}
                  node[actor, left=2pt]{\textit{(Sistema)}} (archivada);

  % --- Rama observaci\'{o}n (salida derecha del rombo) ---
  \draw[->, thick, color=Micolor1!80]
    (dec.east) -- node[etiq, above=2pt]{No --- observa}
                  node[actor, below=2pt]{\textit{(Com.\,Curricular)}} (observada.west);

  % --- Retorno: observada sube y vuelve a En revisi\'{o}n ---
  \draw[->, thick, color=Micolor1!80]
    (observada.north) to[out=90, in=0]
      node[etiq, pos=0.45, right=3pt]{corrige y reenv\'ia}
      node[actor, pos=0.72, right=3pt]{\textit{(Docente)}}
      (revision.east);
\end{tikzpicture}
\shorthandon{<>}
\caption{Diagrama de flujo del ciclo de vida de propuestas en \textsc{macau}.\\
\footnotesize\textit{Fuente: Elaboraci\'{o}n propia.}}
\label{fig:flujo_estados_propuesta}
\end{figure}

\subsection{Campos de una propuesta de asignatura}
\label{subsec:campos_propuesta}

Toda propuesta almacenada en \textsc{macau} se estructura en seis bloques
temáticos que replican la plantilla institucional vigente.

\textbf{Identificación y régimen académico.} Nombre oficial de la asignatura,
carrera, plan de estudios de referencia, año
lectivo, año en la carrera, cuatrimestre, carácter (obligatoria - optativa) y
régimen de cursada (cuatrimestral - anual). Estos atributos ubican la propuesta
dentro de la estructura curricular y permiten verificar su coherencia con el
plan vigente.

\textbf{Carga horaria.} Horas semanales y totales, diferenciadas entre teoría y práctica. Estos valores se utilizan en los controles normativos del módulo de revisión, que verifica la consistencia de la carga horaria declarada con lo establecido en la resolución de la carrera.

\textbf{Equipo docente.} Relación de docentes vinculados con rol en la Asignatura
(Titular, Asociado, Adjunto, Jefe de Trabajos Prácticos, Ayudante). Se almacenan como registros independientes en
\texttt{ProposalTeacher}, lo que permite armar equipos multi-carrera sin
duplicar la propuesta.

\textbf{Fundamentación, Contenidos y Competencias} Se incluyen campos de texto libre para la fundamentación —orientada a la relevancia de la asignatura y su vinculación con el perfil del egresado— y para los contenidos mínimos establecidos en la Resolución de la carrera. Asimismo, se incorporan las competencias genéricas y específicas.

\textbf{Secciones dinámicas.} Tres colecciones almacenadas como JSON
estructurado que modelan el contenido pedagógico variable:
\begin{itemize}[noitemsep, topsep=2pt]
    \item \textit{Unidades temáticas}: identificador, título, descripción de contenidos por unidad, bibliografía básica y complementaria;
    \item \textit{Trabajos prácticos}: Objetivos y Resultados de Aprendizaje a los que aporta, Actividades a desarrollar, Materiales, Ámbito de Práctica;
    \item \textit{Resultados de aprendizaje} (RA): descripción de los resultados de aprendizaje esperados.
\end{itemize}

\textbf{Secciones metodológicas.} Metodología de enseñanza, sistema de
evaluación, bibliografía básica y complementaria, y un campo libre de
observaciones; todos en texto plano editable en la interfaz o importable
desde el documento fuente.


\subsection{Importaci\'{o}n de propuestas existentes}
\label{sec:extraccion}

Uno de los requisitos de adopci\'{o}n cr\'{i}ticos para cualquier plataforma
institucional es la capacidad de incorporar el acervo documental preexistente
sin obligar a la re-ingesta manual. En el contexto de \textsc{macau}, esto
significa que las propuestas de asignatura elaboradas hist\'oricamente con la
plantilla institucional de la universidad o en Google~Docs deben poder migrarse al
sistema conservando la estructura y el contenido original. Para ello se
implement\'{o} un m\'{o}dulo de importaci\'{o}n con dos canales de entrada
que convergen en un pipeline de extracci\'{o}n \'{u}nico.

\subsubsection{Pipeline de extracci\'{o}n desde DOCX}
\label{subsubsec:import_docx}

El canal primario de importaci\'{o}n opera sobre el archivo DOCX generado con
la plantilla institucional. El flujo, iniciado mediante
\texttt{POST /proposals/import-docx}, comprende las siguientes etapas:

\begin{enumerate}[noitemsep, topsep=3pt]
    \item \textbf{Recepci\'{o}n y almacenamiento temporal.} El archivo se
          recibe como \textit{multipart upload}, se persiste en un directorio
          temporal del servidor y se referencia por su ruta absoluta para el
          procesamiento posterior.

    \item \textbf{Detecci\'{o}n de estructura.} El m\'{o}dulo
          \texttt{docx\_import.py} abre el documento con \textit{python-docx}
          y aplica una estrategia de detecci\'{o}n en dos capas: primero
          recorre los p\'{a}rrafos buscando encabezados normalizados que
          coincidan con palabras clave predefinidas
          (\textit{Contenidos M\'{i}nimos}, \textit{Fundamentos},
          \textit{Metodolog\'{i}a}, \textit{Bibliograf\'{i}a}, etc.); luego,
          en los casos en que el contenido se encuentre dentro de tablas en
          lugar de p\'{a}rrafos libres, aplica una segunda b\'{u}squeda sobre
          las celdas de la tabla. Esta dualidad cubre las variantes de formato
          que coexisten en el corpus institucional.

    \item \textbf{Normalizaci\'{o}n.} Para lograr robustez ante variantes
          de escritura, los t\'{e}rminos clave se normalizan eliminando
          acentos, convirtiendo a min\'{u}sculas y colapsando espacios
          m\'{u}ltiples antes de toda comparaci\'{o}n. Esto permite reconocer,
          por ejemplo, \textit{``Metodolog\'{i}a:''}, \textit{``METODOLOGIA''}
          y \textit{``Metodolog\'{i}a de Ense\~{n}anza''} como la misma
          secci\'{o}n objetivo.

    \item \textbf{Mapeo al modelo.} Los bloques de texto extra\'{i}dos se
          asignan a los campos del modelo relacional \texttt{Proposal}:
          secciones de texto libre (\texttt{minimum\_content},
          \texttt{methodology}, \texttt{evaluation}, \texttt{bibliography}),
          colecciones estructuradas en JSON (\texttt{units},
          \texttt{practicals}, \texttt{learning\_outcomes}) y metadatos de
          identificaci\'{o}n (\texttt{career}, \texttt{subject},
          \texttt{year\_of\_career}, \texttt{theoretical\_hours}, etc.).

    \item \textbf{Respuesta de previsualizaci\'{o}n.} El endpoint no persiste
          directamente en la base de datos: devuelve un objeto JSON con los
          datos extra\'{i}dos y un resumen de previsualizaci\'{o}n (conteo de
          unidades, trabajos pr\'{a}cticos, resultados de aprendizaje y
          competencias) para confirmaci\'{o}n por parte del usuario antes del
          guardado definitivo.
\end{enumerate}

\subsubsection{Canal alternativo: importaci\'{o}n desde Google~Docs}
\label{subsubsec:import_gdoc}

El segundo canal, accesible mediante \texttt{POST /proposals/import-gdoc-url},
permite importar documentos alojados en Google~Drive sin necesidad de
descargarlos manualmente. El mecanismo extrae el identificador \'{u}nico del
documento desde la URL recibida (soportando tanto el formato
\texttt{/d/\{id\}/edit} de Google~Docs como el formato
\texttt{?id=\{id\}} de Google~Drive), construye la URL de exportaci\'{o}n
est\'{a}ndar de la API de Google
(\texttt{/export?format=docx}) y descarga el archivo resultante como flujo de
bytes. A partir de ese punto, el contenido binario se canaliza al mismo extractor
\texttt{import\_proposal\_from\_docx} empleado por el canal DOCX, garantizando
un tratamiento homog\'{e}neo de ambas fuentes. Como condici\'{o}n previa, el
documento debe tener visibilidad p\'{u}blica o compartida mediante enlace, dado
que la descarga se realiza sin credenciales OAuth en esta modalidad.

La Figura~\ref{fig:pipeline_importacion} ilustra ambos canales y su
convergencia en el pipeline com\'{u}n de extracci\'{o}n.

\begin{figure}[H]
\centering
\shorthandoff{<>}
\begin{tikzpicture}[
  every node/.style={font=\small},
  fuente/.style={
    rectangle, rounded corners=4pt,
    draw=Micolor1, line width=0.8pt,
    minimum width=3.0cm, minimum height=0.85cm,
    align=center, fill=Micolor1!12
  },
  proceso/.style={
    rectangle, rounded corners=2pt,
    draw=gray!60, line width=0.7pt,
    minimum width=3.2cm, minimum height=0.85cm,
    align=center, fill=gray!8
  },
  salida/.style={
    rectangle, rounded corners=4pt,
    draw=Micolor1, line width=0.8pt,
    minimum width=3.2cm, minimum height=0.85cm,
    align=center, fill=Micolor1!20
  },
  etiq/.style={font=\scriptsize\bfseries},
  actor/.style={font=\scriptsize\itshape, text=gray!70!black}
]
  % --- Fuentes de entrada ---
  \node[fuente] (docx)  at (0, 0)    {Archivo DOCX\\{\scriptsize plantilla institucional}};
  \node[fuente] (gdoc)  at (5.5, 0)  {Google Docs URL\\{\scriptsize documento p\'{u}blico}};

  % --- Paso intermedio GDoc ---
  \node[proceso] (export) at (5.5, -1.8) {Exportar a DOCX\\{\scriptsize Google Drive API}};

  % --- Extractor com\'{u}n ---
  \node[proceso] (extractor) at (2.75, -3.6)
    {Extractor \texttt{docx\_import.py}\\
     {\scriptsize detecci\'{o}n $\cdot$ normalizaci\'{o}n $\cdot$ mapeo}};

  % --- Salida ---
  \node[salida] (preview) at (2.75, -5.4)
    {JSON de previsualizaci\'{o}n\\{\scriptsize + campos modelo \texttt{Proposal}}};

  % --- Flechas ---
  \draw[->, thick, color=Micolor1]
    (docx.south) -- ++(0,-0.6) -| node[etiq, pos=0.25, above]{\texttt{POST /proposals/import-docx}} (extractor.north west);

  \draw[->, thick, color=Micolor1]
    (gdoc.south) -- ++(0,-0.6) -| node[etiq,pos=0.25, above]{\texttt{POST /proposals/import-gdoc-url}} (export.north);

  \draw[->, thick, color=gray!70]
    (export.south) -- ++(0,-0.5) -| (extractor.north east);

  \draw[->, thick, color=Micolor1]
    (extractor) -- node[etiq, right=2pt]{datos extra\'{i}dos} (preview);
\end{tikzpicture}
\shorthandon{<>}
\caption{Pipeline de importaci\'{o}n de propuestas: canales DOCX y Google~Docs con
extractor com\'{u}n. \footnotesize\textit{Fuente: Elaboraci\'{o}n propia.}}
\label{fig:pipeline_importacion}
\end{figure}

La arquitectura de canal \'{u}nico de extracci\'{o}n persigue dos objetivos
complementarios: simplificar el mantenimiento (cualquier mejora al extractor
beneficia autom\'{a}ticamente a ambas v\'{i}as de ingesta) y garantizar que la
calidad del reconocimiento sea homog\'{e}nea independientemente del origen del
documento. La evaluaci\'{o}n cuantitativa del m\'{o}dulo sobre el corpus
documental real se presenta en la Secci\'{o}n~\ref{sec:resultados_extraccion}.

\subsection{Integraci\'{o}n bidireccional con Google Workspace}
\label{subsec:gdocs_sync}

El dise\~{n}o de \textsc{macau} reconoce que, en la pr\'{a}ctica, Google~Docs se ha consolidado como el entorno de edici\'{o}n
colaborativa m\'{a}s difundido. Por ello se implement\'{o} una capa de
integraci\'{o}n bidireccional con Google~Workspace que permite que el documento
en la nube y el registro estructurado en la base de datos coexistan y se
mantengan sincronizados a lo largo del ciclo lectivo, sin obligar al docente a
elegir entre el confort de la edici\'{o}n en l\'{i}nea y los mecanismos de
control y validaci\'{o}n de la plataforma.

\subsubsection{Vinculaci\'{o}n y creaci\'{o}n de documentos}
\label{subsubsec:gdoc_link}

La integraci\'{o}n puede iniciarse en cualquiera de las dos direcciones.
Si el docente dispone de un borrador preexistente en Google~Docs, puede
vincularse mediante el endpoint \texttt{POST /propo\-sals/\{id\}/link-gdoc}:
el sistema extrae el identificador del documento de la URL, descarga y procesa el DOCX
resultante vía la API de exportaci\'{o}n de Google, y verifica que la
asignatura mencionada en el documento coincida con la propuesta en la base
de datos antes de persistir la vinculaci\'{o}n (\texttt{gdoc\_url},
\texttt{gdoc\_hash}, \texttt{gdoc\_last\_synced}). Esta validaci\'{o}n
impide asociar documentos ajenos a la asignatura y deja un registro de
auditor\'{i}a de la operaci\'{o}n.

El flujo inverso se cubre con el endpoint
\texttt{POST /propo\-sals/\{id\}/create-gdoc}: el sistema genera un DOCX a
partir del estado actual de la propuesta en la base de datos, lo sube a la
carpeta de Google~Drive configurada para la carrera (\texttt{DriveSettings})
y convierte el archivo en un Google~Doc editable. El docente obtiene
inmediatamente un enlace de edici\'{o}n en l\'{i}nea con el contenido ya
estructurado. Este flujo requiere credenciales OAuth v\'{a}lidas configuradas
por el administrador (variables \texttt{GOOGLE\_OAUTH\_*}).

\subsubsection{Ciclo de sincronizaci\'{o}n}
\label{subsubsec:gdoc_sync_cycle}

Una vez establecido el v\'{i}nculo, la plataforma expone tres acciones
complementarias para gestionar la divergencia entre ambas copias del
documento:

\begin{description}[noitemsep, topsep=4pt, leftmargin=0pt]
  \item[\texttt{GET /proposals/\{id\}/gdoc-diff} --- Vista previa de cambios
    entrantes.] Descarga el Google~Doc actual, construye un
    \textit{snapshot} estructurado campo a campo y lo compara con el
    \textit{snapshot} local mediante hash de contenido. Retorna un objeto
    \texttt{changes} con los campos que divergen, sus valores actuales y la
    versi\'{o}n del GDoc, e indica qu\'{e} campos requieren revisi\'{o}n
    editorial (por ejemplo, \texttt{minimum\_content} y
    \texttt{teaching\_team}). El usuario puede inspeccionar las diferencias
    antes de comprometerse con el cambio.

  \item[\texttt{POST /proposals/\{id\}/sync-gdoc} --- Ingesta desde el GDoc.]
    Aplica los campos extra\'{i}dos del Google~Doc a la propuesta local,
    actualiza \texttt{gdoc\_hash} y \texttt{gdoc\_last\_synced} y retorna la
    propuesta actualizada. La misma validaci\'{o}n de coincidencia de asignatura
    que en la vinculaci\'{o}n impide sobrescrituras accidentales.

  \item[\texttt{POST /proposals/\{id\}/push-to-gdoc-direct} --- Publicaci\'{o}n
    hacia el GDoc.] El operador selecciona qu\'{e} campos de la plataforma deben
    propagarse al documento en la nube. El sistema genera un DOCX actualizado a
    partir de la plantilla institucional, lo sube a Google~Drive reemplazando el
    contenido del archivo original y actualiza el hash en la base de datos. Esta
    operaci\'{o}n es selectiva: los campos no marcados permanecen con los
    valores presentes en el GDoc.
\end{description}

Adicionalmente, \texttt{gdoc-accept-latest} aplica todos los cambios
detectados en el GDoc de una sola vez (sin revisi\'{o}n campo a campo), y
\texttt{unlink-gdoc} desvincula el documento sin eliminarlo de Drive. Un
endpoint auxiliar, \texttt{gdoc-status}, permite al frontend consultar en
lote el estado de vinculaci\'{o}n de una colecci\'{o}n de propuestas
(\textit{Sincronizado}, \textit{Sin vincular}, \textit{Link perdido},
\textit{Actualizado}, \textit{Cambios sin enviar}).

La Figura~\ref{fig:gdoc_sync} sintetiza los flujos de ambas direcciones de
sincronizaci\'{o}n.

\begin{figure}[H]
\centering
\shorthandoff{<>}
\resizebox{0.25\linewidth}{!}{%
\begin{tikzpicture}[
  x=1cm, y=1cm,
  caja/.style={
    rectangle, rounded corners=4pt,
    minimum width=3.2cm, minimum height=0.9cm,
    align=center, font=\small
  },
  sistema/.style={caja, draw=Micolor1, line width=0.9pt, fill=Micolor1!10},
  externo/.style={caja, draw=gray!55, line width=0.8pt, fill=gray!7},
  proc/.style={caja, draw=Micolor1!45, line width=0.7pt, fill=Micolor1!4},
  flecha/.style={->, thick},
  fbd/.style={flecha, color=Micolor1},
  fgd/.style={flecha, color=gray!60},
  lbl/.style={font=\footnotesize, fill=white, inner sep=1.5pt, align=center},
  sec/.style={font=\small\bfseries, text=Micolor1!75!black},
  lin/.style={gray!30, dashed, line width=0.6pt}
]

%% ─────────────────────────────────────────────────────────
%% Columnas: izq=0  centro=5  der=10
%% Cada sección ocupa al menos 4 cm de alto
%% ─────────────────────────────────────────────────────────

%% SECCIÓN 0 ─ VINCULAR
\node[sec] at (5, 0.55) {0.\enspace Vinculaci\'{o}n};

\node[sistema] (BD0) at (0,   -0.9) {\textsc{macau} BD};
\node[externo] (GD0) at (10,  -0.9) {Google Docs};

% flecha derecha→izquierda (link-gdoc)
\draw[fbd]
  ([yshift= 6pt]GD0.west) -- ([yshift= 6pt]BD0.east)
  node[lbl, midway, above=3pt] {\texttt{link-gdoc}};
\node[lbl, white] at (5, -1.55) {\footnotesize el usuario aporta la URL; el sistema valida};

% flecha izquierda→derecha (create-gdoc)
\draw[fgd]
  ([yshift=-6pt]BD0.east) -- ([yshift=-6pt]GD0.west)
  node[lbl, midway, below=3pt] {\texttt{create-gdoc}};
\node[lbl, white] at (5, -2.4) {\footnotesize genera DOCX en BD, sube a Drive, retorna enlace};

%% línea separadora
\draw[lin] (-1.5, -3.2) -- (11.5, -3.2);

%% SECCIÓN 1 ─ PULL
\node[sec] at (5, -3.75) {1.\enspace PULL \enspace GDoc $\longrightarrow$ BD};

\node[externo] (GD1)  at (10,  -5.1) {Google Docs};
\node[externo] (API1) at (10,  -7.1) {Drive API\\{\scriptsize \texttt{export?format=docx}}};
\node[proc]    (EXT)  at (5,   -7.1) {Extractor\\{\scriptsize \texttt{docx\_import.py}}};
\node[proc]    (DIF)  at (5,   -9.2) {\texttt{gdoc-diff}\\{\scriptsize diff campo a campo}};
\node[sistema] (BD1)  at (0,   -9.2) {\textsc{macau} BD\\{\scriptsize actualizada}};

\draw[fgd] (GD1.south) -- (API1.north)
  node[lbl, right=5pt, midway] {exporta DOCX};

\draw[fgd] (API1.west) -- (EXT.east)
  node[lbl, above=3pt, midway] {DOCX};

\draw[fbd] (EXT.south) -- (DIF.north)
  node[lbl, right=5pt, midway] {snapshot};

\draw[fbd] (DIF.west) -- (BD1.east)
  node[lbl, above=3pt, midway] {\texttt{sync-gdoc}};

% bypass accept-latest
\draw[->, Micolor1!30, thick, dashed]
  (EXT.west) to[out=200, in=30, looseness=0.5]
  node[lbl, pos=0.5, left=2pt] {\scriptsize \texttt{accept-latest}}
  (BD1.north);

%% línea separadora
\draw[lin] (-1.5, -10.6) -- (11.5, -10.6);

%% SECCIÓN 2 ─ PUSH
\node[sec] at (5, -11.15) {2.\enspace PUSH \enspace BD $\longrightarrow$ GDoc};

\node[sistema] (BD2) at (0,   -12.4) {\textsc{macau} BD\\{\scriptsize campos seleccionados}};
\node[proc]    (GEN) at (0,   -14.3) {Generador\\{\scriptsize \texttt{docx\_export.py}}};
\node[externo] (DRV) at (5,   -14.3) {Drive API\\{\scriptsize \texttt{files().update()}}};
\node[externo] (GD2) at (10,  -14.3) {Google Docs\\{\scriptsize contenido actualizado}};

\draw[fbd] (BD2.south) -- (GEN.north)
  node[lbl, right=5pt, midway] {\texttt{push-to-gdoc}};

\draw[fgd] (GEN.east) -- (DRV.west)
  node[lbl, above=3pt, midway] {DOCX};

\draw[fgd] (DRV.east) -- (GD2.west)
  node[lbl, above=3pt, midway] {reemplaza};

\end{tikzpicture}
}% end resizebox
\shorthandon{<>}
\caption{Integraci\'{o}n bidireccional con Google~Workspace: vinculaci\'{o}n inicial,
flujo \textit{pull} (GDoc $\rightarrow$ BD) y flujo \textit{push}
(BD $\rightarrow$ GDoc).
\footnotesize\textit{Fuente: Elaboraci\'{o}n propia.}}
\label{fig:gdoc_sync}
\end{figure}

El estado de la vinculaci\'{o}n se persiste en tres campos del modelo
\texttt{Proposal}:

\begin{description}[noitemsep, topsep=3pt]
  \item[\texttt{gdoc\_url}] URL can\'{o}nica del documento vinculado.
    Se almacena al ejecutar \texttt{link-gdoc} o \texttt{create-gdoc} y es
    la referencia utilizada por todos los endpoints de sincronizaci\'{o}n.

  \item[\texttt{gdoc\_hash}] Hash del contenido del documento en la \'{u}ltima
    sincronizaci\'{o}n exitosa. Permite detectar si el GDoc fue modificado
    desde la \'{u}ltima operaci\'{o}n sin necesidad de descargarlo: si el hash
    actual coincide con el almacenado, no hay cambios pendientes.

  \item[\texttt{gdoc\_status}] Estado de vinculaci\'{o}n visible en la
    interfaz. Puede tomar cinco valores:
    \textit{Sincronizado} (doc vinculado y al d\'{i}a),
    \textit{Sin vincular} (propuesta sin GDoc asociado),
    \textit{Link perdido} (URL registrada pero documento inaccesible),
    \textit{Actualizado} (el GDoc tiene cambios no incorporados a la BD) e
    \textit{Cambios sin enviar} (la BD tiene cambios no propagados al GDoc).
\end{description}

Este esquema permite que el frontend muestre el indicador de estado sin
consultar la API de Google en cada carga de p\'{a}gina: el valor se actualiza
\'{u}nicamente cuando el usuario o el sistema ejecuta una operaci\'{o}n de
sincronizaci\'{o}n expl\'{i}cita.

\subsection{Exportaci\'{o}n de propuestas}
\label{subsec:exportacion}

El ciclo de vida de una propuesta culmina, en t\'{e}rminos documentales,
con la generaci\'{o}n del archivo que se entregar\'{a} a la Comisi\'{o}n
Curricular o se archivar\'{a} en el expediente de acreditaci\'{o}n.
\textsc{macau} expone dos rutas para esta operaci\'{o}n.
La primera, \texttt{GET /proposals/\{id\}/docx}, es un atajo de conveniencia
que delega internamente en la segunda, la ruta gen\'{e}rica
\texttt{GET /proposals/\{id\}/export?format=\textit{fmt}}, la cual soporta
cuatro formatos: \texttt{docx}, \texttt{pdf}, \texttt{json} y \texttt{xml}.
El par\'{a}metro \texttt{format} por defecto es \texttt{docx}, de modo que
ambas rutas son equivalentes cuando no se especifica otro valor.

\medskip
\noindent\textbf{Formatos de documento (DOCX y PDF).}
Para cualquiera de estos dos formatos el sistema carga la plantilla
institucional (reproducida en el Anexo~\ref{anexo:plantilla_docx},
p\'{a}g.~\pageref{anexo:plantilla_docx})\footnote{La plantilla es el archivo
\texttt{Template\,Propuestas.docx} ubicado en el directorio de datos del
servidor; contiene marcadores de la forma \texttt{\{\{campo\}\}} que el
motor reemplaza por los valores de la base de datos.}
e invoca \texttt{generate\_proposal\_\allowbreak docx()} del m\'{o}dulo
\texttt{docx\_export.py}.
La funci\'{o}n recorre todos los bloques del documento (p\'{a}rrafos y
tablas) mediante \texttt{\_iter\_block\_\allowbreak items()}, localiza cada marcador
\texttt{\{\{campo\}\}} con una expresi\'{o}n regular, preserva el formato
del \textit{run} original y lo sustituye por el valor almacenado en la
base de datos.
El resultado es un DOCX que refleja con exactitud el estado en que la propuesta fue
aprobada, descargado con nombre autom\'{a}tico del tipo
\mbox{\texttt{<carrera>\_<asignatura>\_<a\~{n}o>.docx}}.

Si el formato solicitado es PDF, el DOCX generado se convierte mediante
LibreOffice en modo \textit{headless} (\texttt{soffice --headless
--convert-to pdf}); si LibreOffice no est\'{a} disponible, un mecanismo de
\textit{fallback} alternativo (\texttt{convert\_docx\_to\_pdf}) retoma la
conversi\'{o}n.
La cabecera \texttt{X-PDF-Renderer} de la respuesta HTTP informa qu\'{e}
motor se utiliz\'{o}, facilitando el diagn\'{o}stico ante diferencias de
estilo entre entornos de ejecuci\'{o}n.

\medskip
\noindent\textbf{Formatos de datos estructurados (JSON y XML).}
Los formatos \texttt{json} y \texttt{xml} serializan la propuesta completa:
todos los campos, colecciones anidadas (resultados de aprendizaje,
competencias, unidades) y metadatos de estado.
El primero se construye con \texttt{json.dumps(\ldots,\allowbreak~ensure\_ascii=False)},
preservando caracteres Unicode y produciendo salida legible e
indentada.
El segundo emplea \texttt{xml.etree.ElementTree}, recorriendo
recursivamente el dict de la propuesta y convirtiendo cada clave en
un elemento XML v\'{a}lido.
Ninguno de los dos requiere la plantilla DOCX, por lo que son accesibles
incluso si la versi\'{o}n de plantilla no coincide.
Estos formatos est\'{a}n orientados a integraciones con sistemas
institucionales externos (SIU~Guar\'{a}n\'{i}, repositorios de planes de
estudio) sin necesidad de re-extracci\'{o}n desde documentos.

\medskip
La Figura~\ref{fig:exportacion} resume la arquitectura de exportaci\'{o}n:
una \'{u}nica ruta gen\'{e}rica bifurca en cuatro caminos seg\'{u}n el formato
solicitado.

\begin{figure}[H]
\centering
\shorthandoff{<>}
\begin{tikzpicture}[x=1cm, y=1cm,
  cjt/.style={rectangle, rounded corners=4pt, draw=Micolor1, line width=0.9pt,
    minimum width=10.8cm, minimum height=1.05cm,
    align=center, fill=Micolor1!10, font=\small},
  fmt/.style={rectangle, rounded corners=3pt, draw=Micolor1!70, line width=0.8pt,
    minimum width=2.35cm, minimum height=0.75cm,
    align=center, fill=Micolor1!15, font=\small\bfseries},
  mot/.style={rectangle, rounded corners=2pt, draw=gray!50, line width=0.7pt,
    minimum width=2.35cm, minimum height=1.55cm,
    align=center, fill=gray!5, font=\scriptsize},
  sal/.style={rectangle, rounded corners=3pt, draw=Micolor1!35, line width=0.7pt,
    minimum width=2.35cm, minimum height=0.75cm,
    align=center, fill=Micolor1!5, font=\scriptsize\itshape},
  fa/.style={->, gray!60, thick},
  fb/.style={->, Micolor1!55, thick}
]

%% ── Ruta genérica (top) ──────────────────────────────────────────────
\node[cjt] (EP) at (5.4, 0) {
  \texttt{GET /proposals/\{id\}/export\,?\,format=\textit{fmt}}\\[1pt]
  {\scriptsize alias: \texttt{GET /proposals/\{id\}/docx}
    $\;\equiv\;$ \texttt{format=docx}}
};

%% ── Formato (fila 1) ─────────────────────────────────────────────────
\node[fmt] (FD) at (0.5,  -1.9) {docx};
\node[fmt] (FP) at (3.2,  -1.9) {pdf};
\node[fmt] (FJ) at (7.4,  -1.9) {json};
\node[fmt] (FX) at (10.2, -1.9) {xml};

\draw[fb] (EP.south -| FD.north) -- (FD.north);
\draw[fb] (EP.south -| FP.north) -- (FP.north);
\draw[fb] (EP.south -| FJ.north) -- (FJ.north);
\draw[fb] (EP.south -| FX.north) -- (FX.north);

%% ── Motor (fila 2) ───────────────────────────────────────────────────
\node[mot] (MD) at (0.5,  -4.1) {
  \texttt{docx\_export.py}\\[3pt]
  rellena plantilla\\institucional\\campo a campo
};
\node[mot] (MP) at (3.2,  -4.1) {
  \texttt{docx\_export.py}\\[3pt]
  + LibreOffice\\
  \texttt{-{}-headless}\\
  {\tiny cab.~\texttt{X-PDF-Renderer}}
};
\node[mot] (MJ) at (7.4,  -4.1) {
  \texttt{json.dumps}\\[3pt]
  todos los campos\\+ metadatos,\\Unicode preservado
};
\node[mot] (MX) at (10.2, -4.1) {
  \texttt{xml.etree}\\[3pt]
  todos los campos\\+ metadatos,\\UTF-8
};

\draw[fa] (FD.south) -- (MD.north);
\draw[fa] (FP.south) -- (MP.north);
\draw[fa] (FJ.south) -- (MJ.north);
\draw[fa] (FX.south) -- (MX.north);

%% ── Salida (fila 3) ──────────────────────────────────────────────────
\node[sal] (OD) at (0.5,  -5.85) {.docx\\{\tiny revisi\'{o}n curricular}};
\node[sal] (OP) at (3.2,  -5.85) {.pdf\\{\tiny archivo / acreditaci\'{o}n}};
\node[sal] (OJ) at (7.4,  -5.85) {.json\\{\tiny SIU~Guar\'{a}n\'{i} / API}};
\node[sal] (OX) at (10.2, -5.85) {.xml\\{\tiny integraci\'{o}n externa}};

\draw[fa] (MD.south) -- (OD.north);
\draw[fa] (MP.south) -- (OP.north);
\draw[fa] (MJ.south) -- (OJ.north);
\draw[fa] (MX.south) -- (OX.north);

\end{tikzpicture}
\shorthandon{<>}
\caption{Arquitectura de exportaci\'{o}n: la ruta gen\'{e}rica
  \texttt{/export?format=} bifurca en cuatro caminos seg\'{u}n el formato
  solicitado.
  \footnotesize\textit{Fuente: Elaboraci\'{o}n propia.}}
\label{fig:exportacion}
\end{figure}

La simetr\'{i}a entre importaci\'{o}n y exportaci\'{o}n --- ambas operando
sobre la misma plantilla y el mismo modelo relacional --- asegura que una
propuesta importada desde un DOCX institucional y luego exportada de vuelta
a DOCX produzca un documento estructuralmente equivalente al original, con la
\'{u}nica diferencia del contenido que haya sido editado o enriquecido por la
plataforma.

%% =============================================================================
\section{Controles normativos e inteligentes sobre propuestas de Asignatura}
\label{sec:controles}
%% =============================================================================

El proceso de revisi\'on en \textsc{macau} combina dos capas complementarias.
La primera capa define \textbf{controles normativos fijos}, derivados de los
crite\-rios de consistencia curricular y los est\'andares de acreditaci\'on
exigidos por la ~\cite{ordenanza75} para las carreras de ingenier\'ia. La segunda capa incorpora \textbf{controles
inteligentes configurables}, cuyo comportamiento se ajusta por carrera mediante
prompts administrables por perfiles directivos. En conjunto, este esquema evita
que la validación dependa exclusivamente del criterio individual del revisor y
produce un resultado comparable entre asignaturas y entre ciclos lectivos.

Operativamente, ambos mecanismos convergen en un flujo único de ejecución y
registro: el backend corre la batería de controles, persiste cada resultado por
propuesta y control en \texttt{ProposalIntelligentControlResult}, y consolida la
salida en un resumen tipado
(\texttt{ProposalIntelligentControlsSummary}) consumido por la interfaz.
El panel de revisión, por lo tanto, no solo indica si una propuesta cumple o
no cumple, sino que preserva evidencia textual de qué falló, por qué falló y
qué corrección se sugirió en cada caso. Las capturas de interfaz se incluyen en
el Anexo~\ref{anexo:capturas}.

\subsection{Controles normativos fijos}
\label{subsec:controles_normativos}

Los controles normativos fijos implementan validaciones estructurales de base,
con especial foco en completitud, trazabilidad y coherencia interna. Entre las
comprobaciones más relevantes se incluyen: presencia de secciones obligatorias,
consistencia mínima de carga horaria, equipo docente formado al menos por un Profesor, existencia de resultados de
aprendizaje, y alineación entre competencias declaradas y estructura de la
propuesta. Aunque la ejecución técnica comparte la misma infraestructura de los
controles inteligentes, su propósito metodológico es distinto: establecen el
\textit{piso de cumplimiento} que toda propuesta debe satisfacer antes de
avanzar a revisión cualitativa.

En términos de API, el ciclo de revisión expone tres operaciones principales:
\begin{itemize}[noitemsep, topsep=2pt]
    \item \texttt{POST .../run}: lanza la corrida de controles activos sobre
      la propuesta;
    \item \texttt{GET .../results}: recupera el estado consolidado de todos
      los resultados;
    \item \texttt{PATCH .../results/\{result\_id\}}: permite la revisión manual
      de un resultado puntual por parte de perfiles con atribuciones de
      supervisión.
\end{itemize}
donde \texttt{.../} abrevia el prefijo
\texttt{/proposals/}\allowbreak\texttt{\{id\}/intelligent-controls}.

Esta última operación no elimina trazabilidad: al actualizar un resultado, el
sistema conserva el registro histórico y actualiza la marca temporal de
verificación, lo que permite auditar excepciones justificadas sin perder el
comportamiento de base del esquema normativo.

\subsection{Controles inteligentes configurables}
\label{subsec:controles_inteligentes}

Los controles inteligentes configurables se modelan con la entidad
\texttt{IntelligentControl}. Cada instancia define un tópico temático
(\textit{p.\,ej.}\ competencias, bibliografía, equipo docente), un nombre
operativo, una instrucción en lenguaje natural que actúa como
\textit{prompt de control}, y parámetros de orden y activación.
Esta separación entre lógica de negocio y código fuente permite que
la carrera evolucione sus criterios de revisión sin modificar la
arquitectura del \textit{backend}.

La gestión del catálogo de controles se resuelve mediante endpoints CRUD
explícitos:
\begin{itemize}[noitemsep, topsep=2pt]
    \item \texttt{GET} y \texttt{POST /intelligent-controls}: listado y alta;
    \item \texttt{PATCH} y \texttt{DELETE /intelligent-controls/\{id\}}:
      actualización y baja;
    \item \texttt{GET} y \texttt{PATCH /intelligent-controls/settings}:
      lectura y escritura de la configuración global de ejecución.
\end{itemize}

\medskip
\noindent\textbf{Integración con el modelo de lenguaje.}

En tiempo de ejecución, el endpoint
\texttt{POST .../intelligent-controls/run} itera sobre los controles activos
de la propuesta, construye el contexto relevante por tópico mediante
\texttt{build\_topic\_payload()} e invoca la función
\texttt{evaluate\_control\_\allowbreak with\_llm()}, que se comunica con la
API de OpenAI a través del cliente configurado con la clave de entorno
\texttt{OPENAI\_API\_KEY}.

Para cada control se ensamblan dos mensajes:
un \textit{system prompt} que instruye al modelo a actuar como evaluador
académico estricto y a devolver exclusivamente JSON con los campos
\texttt{pass}, \texttt{what\_failed}, \texttt{why\_failed},
\texttt{suggestion}, \texttt{proposed\_text} y \texttt{summary};
y un \textit{user prompt} que incluye la regla del control, los datos
del tópico serializado en JSON y, cuando corresponde, contexto asociado
de otros tópicos.
La respuesta se parsea y se persiste en
\texttt{ProposalIntelligentControlResult}, preservando tanto el veredicto
binario (\texttt{pass}) como la justificación textual completa.

\medskip
\noindent\textbf{Modos de ejecución.}
La configuración global (\texttt{IntelligentControlSettings}) define tres
perfiles de ejecución, denominados con nombres de animales como convención
interna del sistema:

\begin{description}[noitemsep, leftmargin=1.4cm, labelwidth=1.2cm]
  \item[\textit{Guepardo}] Perfil rápido y de bajo costo.
    Usa por defecto \texttt{gpt-4o-mini} con temperatura~0{,}15 y límite de
    420~\textit{tokens}.
    Es el modo por defecto asignado a docentes.
  \item[\textit{Delfín}]  Perfil equilibrado entre velocidad y profundidad.
    Usa por defecto \texttt{gpt-4o-mini} con temperatura~0{,}10 y límite de
    500~\textit{tokens}.
    Es el modo por defecto asignado a directivos.
  \item[\textit{Ballena}] Perfil de máxima profundidad de análisis.
    Usa por defecto \texttt{gpt-4o} con temperatura~0{,}10 y límite de
    700~\textit{tokens}.

\end{description}

Cada modo expone su propio modelo, temperatura y límite de tokens como
parámetros configurables de forma independiente en
\texttt{IntelligentControlSettings}, lo que permite ajustar la relación
costo\,/\,calidad según el perfil de usuario activo sin intervención en el
código.
Si el modo indicado no es válido, el sistema normaliza a \textit{Delfín}
como valor de retorno seguro.

El valor agregado de esta capa no reside en una etiqueta binaria, sino en
la devolución estructurada: cada resultado incluye estado de aprobación,
descripción del incumplimiento, justificación textual y sugerencia de
mejora. De este modo, el control inteligente cumple simultáneamente dos
funciones: verificación técnica y asistencia de redacción orientada a
norma.

%% =============================================================================
\section{Arquitectura de agentes (MAS)}
\label{sec:agentes}
%% =============================================================================

La plataforma organiza las capacidades de inteligencia artificial en cuatro
agentes especializados que operan de forma coordinada bajo el backend FastAPI.
Cada agente encapsula una responsabilidad bien delimitada y expone su
funcionalidad a través de la API REST; el frontend actúa como orquestador
de la interacción del usuario, invocando al agente adecuado según la acción
solicitada.
La Figura~\ref{fig:mas_agentes} muestra la arquitectura global del sistema
multi-agente.

\shorthandoff{<>}
\begin{figure}[ht]
\centering
\resizebox{\linewidth}{!}{%
\begin{tikzpicture}[
  x=1cm, y=1cm,
  font=\normalsize,
  box/.style={draw, rounded corners=5pt, minimum width=3.8cm,
              minimum height=1.1cm, align=center, fill=gray!10, line width=0.6pt},
  agent/.style={draw, rounded corners=5pt, minimum width=3.2cm,
                minimum height=1.8cm, align=center, text width=2.8cm,
                fill=Micolor1!12, draw=Micolor1, line width=0.8pt},
  api/.style={draw, rounded corners=5pt, minimum width=3.2cm,
              minimum height=1.1cm, align=center,
              fill=blue!8, draw=blue!60, line width=0.6pt},
  arr/.style={->, >=stealth, line width=1.2pt},
  dasharr/.style={<->, >=stealth, dashed, draw=blue!60, line width=1.0pt},
  lbl/.style={font=\small, inner sep=3pt, fill=white}
]

%% ——— Capa 1: Usuario ———
\node[box] (user) at (8.0, 10.0) {Usuario};

%% ——— Capa 2: Frontend ———
\node[box, minimum width=5.0cm] (fe) at (8.0, 8.2) {Frontend\\(React / Vite)};

%% ——— Capa 3: Backend + OpenAI ———
\node[box, minimum width=7.0cm] (be)  at (8.0,  6.2) {Backend FastAPI\\(MCP Host)};
\node[api, minimum width=3.6cm] (oai) at (16.5, 6.2) {OpenAI API};

%% ——— Capa 4: Cuatro agentes en fila (más separados) ———
\node[agent] (aext) at ( 1.5, 3.2) {Agente\\de\\Extracción};
\node[agent] (allm) at ( 6.0, 3.2) {Agente\\de\\Asistencia LLM};
\node[agent] (asem) at (10.5, 3.2) {Agente\\de\\Búsqueda\\Semántica};
\node[agent] (anot) at (15.0, 3.2) {Agente\\de\\Notificación};

%% ——— Capa 5: Gmail ———
\node[api] (gmail) at (15.0, 1.0) {Gmail API};

%% ——— Flechas verticales principales ———
\draw[arr] (user.south) -- node[lbl, right]{HTTP} (fe.north);
\draw[arr] (fe.south)   -- node[lbl, right]{REST} (be.north);

%% ——— Backend ↔ OpenAI ———
\draw[dasharr] (be.east) -- node[lbl, above]{Llamadas LLM} (oai.west);

%% ——— Backend → agentes ———
\draw[arr] (be.south) -- ++(0,-0.6) -| (aext.north);
\draw[arr] (be.south) -- ++(0,-0.6) -| (allm.north);
\draw[arr] (be.south) -- ++(0,-0.6) -| (asem.north);
\draw[arr] (be.south) -- ++(0,-0.6) -| (anot.north);

%% ——— Notificación → Gmail ———
\draw[arr] (anot.south) -- (gmail.north);

\end{tikzpicture}%
}
\caption{Arquitectura del sistema multiagente. \footnotesize\textit{Fuente: Elaboraci\'{o}n propia.}}
\label{fig:mas_agentes}
\end{figure}
\shorthandon{<>}

\subsection{Agente de Extracción}
\label{subsec:agente_extraccion}

El Agente de Extracción, implementado en \texttt{agents/extract.py}, es
responsable de transformar documentos institucionales en datos estructurados
listos para ser almacenados en la base de datos.
Su función principal es \texttt{process\_file()}: recibe la ruta de un archivo
DOCX importado y lo procesa mediante el módulo \texttt{docx\_import}, que
aplica las reglas de extracción definidas para el formato de plantilla
institucional (identificación de secciones, tablas de trabajos prácticos,
bloques de contenidos mínimos y equipo docente).
Los prácticos extraídos se vinculan directamente a la propuesta
correspondiente actualizando la columna JSON \texttt{practicals} del modelo
\texttt{Proposal}.

La invocación del agente se integra en el flujo de importación de la API:
cuando el endpoint \texttt{POST /proposals/import-docx} recibe un archivo,
lanza \texttt{extract\_agent.process\_file()} como tarea en segundo plano
(\texttt{BackgroundTasks} de FastAPI), de modo que la respuesta HTTP al
cliente no bloquea el procesamiento.
Este diseño desacopla la carga del documento del tiempo de respuesta
percibido por el usuario.

El agente también expone \texttt{query\_local()}, interfaz de búsqueda sobre
la base de conocimiento local que el Agente de Búsqueda Semántica y el de
Asistencia LLM invocan para recuperar evidencia contextual relevante antes de
construir sus prompts (detalle en §\ref{sec:rag}).

\subsection{Agente de Asistencia LLM}
\label{subsec:agente_llm}

El Agente de Asistencia LLM provee tres funciones de redacción asistida que
el usuario puede activar en cualquier campo de texto de la propuesta.
En todos los casos, el agente inyecta el contexto institucional disponible
en el \emph{system prompt}: nombre de la carrera, perfil de egreso, lista
de competencias del plan vigente y política de terminología institucional
(\texttt{TERMINOLOGY\_POLICY\_PROMPT}).
Este contexto garantiza que las respuestas generadas estén alineadas con los
estándares CONEAU y el modelo educativo de la institución, no con respuestas
de propósito general.

\begin{description}[noitemsep, leftmargin=1.6cm, labelwidth=1.4cm]
  \item[Generación (\texttt{/ai-generate})]
    Genera el contenido de un campo desde cero a partir de una instrucción
    del usuario.
    Es el único modo aplicable cuando el campo está vacío: los tres modos
    restantes requieren texto previo sobre el que operar.
    El docente describe su intención (por ejemplo:
    \textit{«redactá los objetivos para una asignatura de estructuras de
    datos de primer año de Ingeniería Mecatrónica»}) y el agente devuelve
    el texto listo para insertar o ajustar.

  \item[Corrección ortográfica]
    Corrige únicamente los errores ortográficos del texto existente sin
    alterar el estilo, el tono ni la estructura de las oraciones.
    Es el modo menos invasivo: el contenido y la redacción original se
    conservan íntegramente.

  \item[Corrección gramatical]
    Corrige errores gramaticales y de puntuación y mejora la cohesión
    cuando es necesario, respetando el sentido original.
    A diferencia del modo ortográfico, puede reorganizar levemente las
    oraciones para ganar fluidez.

  \item[Reformulación (\texttt{/ai-reformulate})]
    Reescribe el texto de forma notablemente más clara, profesional y
    fluida, preservando el contenido semántico del original.
    Es el modo más transformador: puede estructurar las ideas en párrafos
    diferenciados e introducir listas cuando corresponde.
    La política de terminología evita que el modelo introduzca expresiones
    ajenas al vocabulario institucional.
\end{description}

Un quinto endpoint, \texttt{POST /notifications/email/ai-assist}, aplica
los mismos tres modos operativos sobre texto (ortografía, gramática y
reformulación) al módulo de notificaciones: permite al director corregir
o reformular el borrador de un correo institucional y recibir el resultado
en HTML listo para enviar por Gmail.
La respuesta incluye marcado semántico (\texttt{<strong>}, \texttt{<em>},
\texttt{<p>}) para que el cliente de correo lo renderice correctamente.

\subsection{Agente de Notificación}
\label{subsec:agente_notificacion}

El Agente de Notificación integra la plataforma con la cuenta de Gmail del
director mediante el flujo OAuth 2.0 estándar de Google.
Una vez autorizado, el backend almacena el token de acceso y lo renueva
automáticamente; el endpoint \texttt{GET /notifications/gmail/status}
permite verificar el estado de la conexión en cualquier momento.

El envío se realiza a través de \texttt{POST /notifications/gmail/send},
que acepta destinatario, asunto y cuerpo HTML.
Los casos de uso principales son: docentes con propuestas incompletas fuera
de término, instrumentos evaluativos pendientes de entrega, y actividades
del plan de trabajo de autoevaluación próximas a vencer.
En todos estos escenarios el sistema puede generar el contenido del correo
mediante el Agente de Asistencia LLM (\texttt{/notifications/email/ai-assist})
antes de enviarlo, produciendo mensajes contextualizados con el nombre del
destinatario, la asignatura involucrada y la lista de ítems faltantes, sin
intervención de redacción manual por parte del director.

\subsection{Agente de Búsqueda Semántica}
\label{subsec:agente_busqueda}

La búsqueda semántica se expone a través del endpoint
\texttt{POST /search} y constituye la interfaz unificada para consultas
en lenguaje natural sobre todo el corpus institucional indexado.
A diferencia de una búsqueda por palabras clave, el agente comprende el
\emph{significado} de la consulta y devuelve resultados aunque no haya
coincidencia lexical exacta.

Las consultas típicas incluyen:
\begin{itemize}[noitemsep]
  \item \textit{«¿En qué asignaturas se trabaja aprendizaje basado en
    proyectos?»} — el agente cruza propuestas indexadas y devuelve los
    fragmentos de metodología o fundamentación que contienen ese concepto,
    indicando asignatura, año y cuatrimestre.
  \item \textit{«Mostrame todo lo que involucra al docente García»} —
    recupera propuestas donde ese docente figura como responsable o adjunto,
    proyectos de investigación vinculados e instrumentos evaluativos
    asociados.
  \item \textit{«Estándares CONEAU relacionados con formación práctica»} —
    recupera fragmentos del documento de estándares subidos al módulo de
    acreditación, junto con las evidencias institucionales que los
    satisfacen.
\end{itemize}

El pipeline completo de indexación y consulta se describe en
§\ref{sec:rag}; en esta capa de agente el resultado recuperado puede
inyectarse como contexto al Agente de Asistencia LLM, cerrando el ciclo
de generación aumentada por recuperación (RAG).

%% =============================================================================
\section{Integración MCP — capa de interoperabilidad}
\label{sec:mcp_implementacion}
%% =============================================================================

El \emph{Model Context Protocol} (MCP) es el mecanismo que permite
al modelo de lenguaje operar sobre datos institucionales de forma estructurada
y auditable~\cite{hou2025model}.
El backend FastAPI actúa como \textbf{MCP Host}: instancia procesos Python
independientes (servidores MCP) y les envía mensajes JSON-RPC 2.0 por
\emph{stdio} cuando el modelo necesita leer o escribir datos.
Cada servidor encapsula un dominio específico y accede a la base de datos
SQLite; el modelo nunca escribe directamente sobre los datos, sino que invoca
herramientas (\emph{Tools}) con parámetros validados.
La Figura~\ref{fig:mcp_topologia} muestra la topología resultante.

\shorthandoff{<>}
\begin{figure}[H]
\centering
\resizebox{0.8\linewidth}{!}{%
\begin{tikzpicture}[
  x=1cm, y=1cm,
  font=\normalsize,
  host/.style={draw, rounded corners=5pt, minimum width=5.2cm,
               minimum height=1.1cm, align=center, fill=gray!10, line width=0.7pt},
  srv/.style={draw, rounded corners=5pt, minimum width=4.6cm,
              minimum height=1.1cm, align=center,
              fill=Micolor1!12, draw=Micolor1, line width=0.8pt},
  arr/.style={->, >=stealth, line width=1.2pt},
  dasharr/.style={->, >=stealth, dashed, gray, line width=0.9pt}
]

%% ——— Fila 1: Backend (Host) ———
\node[host] (host) at (9.0, 7.0) {Backend FastAPI (MCP Host)};

%% ——— Fila 2: Tres servidores MCP ———
\node[srv] (sprop) at (2.0, 4.5) {Server Propuestas};
\node[srv] (snorm) at (9.0, 4.5) {Server Normativa};
\node[srv] (saud)  at (16.0, 4.5) {Server Auditoría};

%% Host → Servers
\draw[arr] (host.south) -- ++(0,-0.6) -| (sprop.north);
\draw[arr] (host.south) -- ++(0,-0.6)  -- (snorm.north);
\draw[arr] (host.south) -- ++(0,-0.6) -| (saud.north);

%% ——— Fila 3: Base de datos ———
\node[host, minimum width=5.4cm] (db) at (9.0, 2.0) {Base de datos (SQLite)};

\draw[dasharr] (sprop.south) |- (db.west);
\draw[dasharr] (snorm.south) -- (db.north);
\draw[dasharr] (saud.south)  |- (db.east);

\end{tikzpicture}%
}
\caption{Topología MCP de la plataforma. \textit{Fuente: elaboración propia.}}
\label{fig:mcp_topologia}
\end{figure}
\shorthandon{<>}

\subsection{Servidor de Propuestas}
\label{subsec:mcp_propuestas}

Expone las operaciones sobre el modelo \texttt{Proposal}.
Sus Tools principales son \texttt{leer\_propuesta} (devuelve todos los campos
de una propuesta dado su \texttt{id}) y \texttt{actualizar\_propuesta}
(modifica campos específicos sin sobrescribir el resto), lo que permite al
modelo aplicar sugerencias de redacción de forma quirúrgica sin perder
el estado previo del documento.

\subsection{Servidor de Normativa}
\label{subsec:mcp_normativa}

Indexa los documentos regulatorios vigentes: Ordenanza 75/23, estándares
CONEAU para Ingeniería en Informática y catálogo de competencias CONFEDI.
La Tool \texttt{buscar\_estandar} realiza búsqueda semántica sobre ese corpus
y devuelve los fragmentos más relevantes para una consulta dada.
El corpus puede actualizarse sin reiniciar el backend.

\subsection{Servidor de Auditoría}
\label{subsec:mcp_auditoria}

Registra toda acción que involucra al modelo de lenguaje.
La Tool \texttt{registrar\_evento} escribe un log inmutable con: agente
invocante, acción, \texttt{id} de propuesta, usuario, marca temporal UTC y
\emph{hash} del payload.
La Tool \texttt{historial\_propuesta} devuelve la secuencia de eventos de una
propuesta, permitiendo reconstruir qué generó el sistema, cuándo y con qué
parámetros~\cite{iso_iec_42001_2023}.

\subsection{Transporte y protocolo}
\label{subsec:mcp_transporte}

En desarrollo, cada servidor MCP corre como proceso Python independiente
instanciado vía \texttt{subprocess}; el intercambio es JSON-RPC 2.0 por
entrada/salida estándar (\emph{stdio}).
Para despliegue distribuido el protocolo soporta HTTP+SSE, permitiendo
ejecutar los servidores en instancias separadas sin modificar la lógica
de los agentes.



%% =============================================================================
\section{Módulo RAG — búsqueda semántica sobre documentación}
\label{sec:rag}
%% =============================================================================

La búsqueda por palabras clave devuelve resultados solo cuando la consulta
coincide lexicalmente con el texto almacenado.
Un usuario que busca \textit{«resolución de problemas en contexto real»}
no obtiene resultados si las propuestas hablan de \textit{«situaciones
problemáticas auténticas»}.
El módulo de \emph{Retrieval-Augmented Generation} (RAG) soluciona esta
limitación representando tanto las consultas como los documentos como
vectores en un espacio semántico de alta dimensión, y midiendo la similitud
de coseno entre ellos para encontrar los fragmentos más relevantes
independientemente de la formulación exacta.

\subsection{Corpus indexado}
\label{subsec:rag_corpus}

El sistema mantiene un índice vectorial sobre cuatro tipos de contenido
institucional:

\begin{description}[noitemsep, leftmargin=1.6cm, labelwidth=1.4cm]
  \item[Propuestas] Todos los campos de texto libre de las propuestas de
    asignatura almacenadas: fundamentación, objetivos, contenidos mínimos,
    metodología y bibliografía.
    Cada fragmento conserva metadatos de asignatura, carrera, año y
    cuatrimestre, lo que permite filtrar los resultados por dimensión
    institucional.

  \item[Normativa] Documentos regulatorios subidos al módulo de
    acreditación: Ordenanza 75/23, resoluciones ministeriales,
    estándares CONEAU, criterios de evaluación CONFEDI.
    El índice permite responder consultas como \textit{«¿qué dice la
    normativa sobre carga horaría de
     práctica mínima?»} sin abrir el documento original.

  \item[Competencias] Las descripciones textuales del catálogo de
    competencias (genéricas y específicas por carrera) se indexan para
    que el sistema pueda sugerir qué competencias trabaja una propuesta
    a partir de su contenido, o qué propuestas desarrollan una
    competencia específica.

  \item[Evidencias de acreditación] Los documentos subidos al registro
    de evidencias (informes institucionales, memorias de gestión,
    certificaciones) se incorporan al índice para que una búsqueda
    pueda cruzar propuestas con la documentación de respaldo que las
    sustenta.
\end{description}

\subsection{Pipeline de indexación}
\label{subsec:rag_indexacion}

La Figura~\ref{fig:rag_pipeline} ilustra el flujo completo desde el
documento fuente hasta el índice vectorial y la respuesta al usuario.

\shorthandoff{<>}
\begin{figure}[H]
\centering
\resizebox{\linewidth}{!}{%
\begin{tikzpicture}[
  x=1cm, y=1cm,
  font=\normalsize,
  pstep/.style={draw, rounded corners=4pt, minimum width=3.8cm,
                minimum height=1.1cm, align=center, fill=gray!10, line width=0.7pt},
  rag/.style={draw, rounded corners=4pt, minimum width=3.8cm,
              minimum height=1.1cm, align=center,
              fill=Micolor1!12, draw=Micolor1, line width=0.8pt},
  arr/.style={->, >=stealth, line width=1.2pt},
  lbl/.style={font=\small\itshape, fill=white, inner sep=2pt}
]

%% ---- FILA SUPERIOR: indexación (offline) ----
\node[font=\small\bfseries] at (-0.5, 5.8) {Indexación (offline)};
\node[pstep] (doc)   at ( 1.5, 5.0) {Documento\\fuente};
\node[pstep] (chunk) at ( 6.5, 5.0) {Chunking\\(512 tok., solapam. 64)};
\node[rag]   (embed) at (11.5, 5.0) {Embedding\\text-embedding-3-small};
\node[rag]   (store) at (16.5, 5.0) {Vector Store\\(HNSW, $d=1536$)};

\draw[arr] (doc)   -- (chunk);
\draw[arr] (chunk) -- (embed);
\draw[arr] (embed) -- (store);

%% ---- FILA INFERIOR: consulta (online) ----
\node[font=\small\bfseries] at (-0.5, 2.2) {Consulta (online)};
\node[pstep] (query)  at ( 1.5, 1.5) {Consulta\\del usuario};
\node[rag]   (qembed) at ( 6.5, 1.5) {Embedding\\consulta};
\node[rag]   (ann)    at (11.5, 1.5) {ANN Search\\(top-$k=5$)};
\node[pstep] (frags)  at (16.5, 1.5) {Fragmentos\\recuperados};
\node[pstep] (llm)    at (16.5,-0.5) {LLM + contexto\\(respuesta)};

\draw[arr] (query)  -- (qembed);
\draw[arr] (qembed) -- (ann);
\draw[arr] (ann)    -- node[lbl, above]{similitud coseno} (frags);
\draw[arr] (frags)  -- (llm);

%% ---- Conexión índice → búsqueda ----
\draw[arr, dashed, gray, line width=0.9pt]
  (store.south) -- ++(0,-0.8) -| node[lbl, below right]{índice vectorial} (ann.north);

\end{tikzpicture}%
}
\caption{Pipeline RAG de la plataforma. \textit{Fuente: elaboración propia.}}
\label{fig:rag_pipeline}
\end{figure}
\shorthandon{<>}

El proceso de indexación opera en segundo plano cada vez que se incorpora
un nuevo documento.
El texto completo se divide en fragmentos de 512~\emph{tokens} con
superposición de 64~\emph{tokens} para no perder contexto en los bordes
de los cortes.
Cada fragmento se convierte en un vector de 1536 dimensiones con el modelo
\texttt{text-embedding-3-small} de OpenAI y se inserta en el índice
HNSW (Hierarchical Navigable Small World), una estructura de datos diseñada
para búsqueda aproximada de vecinos más cercanos con complejidad
$O(\log n)$ en consulta~\cite{malkov2018efficient}.
Junto al vector se almacenan los metadatos del fragmento: tipo de corpus
(propuesta / normativa / competencia / evidencia), identificador del
documento origen, posición dentro del documento y texto plano del fragmento.

\subsection{Pipeline de consulta y casos de uso}
\label{subsec:rag_consulta}

Cuando el usuario envía una consulta al endpoint \texttt{POST /search},
el sistema aplica el mismo modelo de embedding sobre el texto de la
consulta y ejecuta una búsqueda ANN sobre el índice, recuperando los
cinco fragmentos con mayor similitud de coseno.
Los resultados se ordenan por score e incluyen el fragmento de texto,
el documento de origen y los metadatos asociados.


Cuando la consulta requiere una respuesta elaborada y no solo la lista de
fragmentos (por ejemplo, \textit{«redactá un resumen de cómo se trabaja la
competencia CG1 a lo largo del plan de estudios»}), los fragmentos
recuperados se inyectan como contexto en el prompt del Agente de Asistencia
LLM, cerrando el ciclo RAG: el modelo genera texto fundamentado en evidencia
real del corpus institucional, no en conocimiento genérico externo.

%% =============================================================================
\section{Gestión del plan de estudios y competencias}
\label{sec:plan_estudios}
%% =============================================================================

Los procesos de acreditación ante la CONEAU exigen que la institución
demuestre, con evidencia documental, que el perfil de egreso declarado
en el plan de estudios es efectivamente alcanzado por el conjunto de
asignaturas que lo conforman.
Sin embargo, en el contexto observado ese análisis se realizaba de forma
reactiva: solo cuando se iniciaba un ciclo de acreditación se intentaba
reconstruir la estructura curricular a partir de documentos dispersos
—planillas Excel sin control de versiones, actas de Comisión Curricular
y propuestas de cátedra no correlacionadas entre sí.
Este módulo centraliza esa información y la hace consultable en tiempo
real, constituyendo el eje estructural sobre el que se apoyan los demás
módulos de la plataforma.

\subsection{Modelo de datos y estructura jerárquica}
\label{subsec:modelo_plan}

El plan de estudios se representa mediante una jerarquía de cuatro
entidades relacionadas en cascada:

\begin{enumerate}
  \item \textbf{\texttt{StudyPlan}}: unidad raíz que identifica la carrera,
  la resolución de aprobación y el año de vigencia.
  \item \textbf{\texttt{StudyYear}}: año académico dentro del plan
  (1.º, 2.º, \ldots, N años según la carrera).
  \item \textbf{\texttt{StudyTerm}}: período lectivo dentro del año
  (cuatrimestre, trimestre o anual según la modalidad institucional).
  \item \textbf{\texttt{StudySubject}}: asignatura con carga horaria
  teórica, práctica y de laboratorio expresadas en horas semanales.
  Las correlatividades entre asignaturas se registran en la entidad
  \texttt{StudySubjectPrerequisite}, que modela la dependencia curricular
  como un grafo dirigido acíclico.
\end{enumerate}

Esta estructura permite representar fielmente cualquier plan vigente en
la institución sin imponer restricciones sobre la cantidad de años,
períodos o asignaturas por período.

\subsection{Gestión de versiones del plan}
\label{subsec:versiones_plan}

Un aspecto crítico en contextos de acreditación es la necesidad de
mantener simultáneamente el plan vigente y las versiones anteriores,
dado que las propuestas de cátedra aprobadas en un ciclo permanecen
asociadas al plan que les dio origen.
La entidad \texttt{StudyPlanStorage} implementa un mecanismo de
versionado: cada carga genera una nueva instancia del plan completo
que queda almacenada con su marca temporal, y la operación
\texttt{POST /study-plans-storage/\{id\}/activate} designa cuál
es la versión vigente sin eliminar las anteriores.
Esto permite comparar planes entre ciclos de acreditación,
identificar qué asignaturas se incorporaron o modificaron y mantener
la trazabilidad histórica que exige la documentación CONEAU.

La vinculación entre una propuesta de cátedra existente y la asignatura
correspondiente en el plan activo se realiza mediante
\texttt{POST /proposals/\{id\}/sync-to-study-plan}, operación que
actualiza la referencia sin modificar el contenido de la propuesta.

\subsection{API del módulo}
\label{subsec:api_plan}

La Tabla~\ref{tab:api_plan} detalla los endpoints del módulo de plan
de estudios, organizados por sub-recurso.

\begin{table}[ht]
\centering
\caption{Endpoints del módulo de plan de estudios.}
\label{tab:api_plan}
\small
\begin{tabular}{llp{5.8cm}}
\toprule
\textbf{Método} & \textbf{Ruta} & \textbf{Descripción} \\
\midrule
\texttt{GET}    & \texttt{/study-plans}                              & Lista todos los planes de la carrera \\
\texttt{POST}   & \texttt{/study-plans}                              & Crea un nuevo plan \\
\texttt{GET}    & \texttt{/study-plans/\{id\}}                       & Recupera el plan completo con su jerarquía \\
\texttt{POST}   & \texttt{/study-plans/\{id\}/years}                 & Agrega un año al plan \\
\texttt{POST}   & \texttt{/study-years/\{id\}/terms}                 & Agrega un período a un año \\
\texttt{POST}   & \texttt{/study-terms/\{id\}/subjects}              & Agrega una asignatura a un período \\
\texttt{PATCH}  & \texttt{/study-subjects/\{id\}}                    & Modifica datos de una asignatura \\
\texttt{DELETE} & \texttt{/study-subjects/\{id\}}                    & Elimina una asignatura \\
\texttt{GET}    & \texttt{/study-subjects}                           & Lista asignaturas con filtros \\
\texttt{GET}    & \texttt{/study-plans-storage}                      & Lista versiones almacenadas \\
\texttt{POST}   & \texttt{/study-plans-storage}                      & Guarda una nueva versión \\
\texttt{POST}   & \texttt{/study-plans-storage/\{id\}/activate}      & Designa la versión vigente \\
\texttt{DELETE} & \texttt{/study-plans-storage/\{id\}}               & Elimina una versión archivada \\
\texttt{POST}   & \texttt{/proposals/\{id\}/sync-to-study-plan}      & Vincula propuesta con asignatura del plan activo \\
\bottomrule
\end{tabular}
\end{table}

El endpoint \texttt{GET /study-plans/\{id\}} devuelve el árbol completo
del plan (años, períodos y asignaturas con correlatividades) en una
única respuesta, lo que evita que el frontend realice múltiples llamadas
para construir la vista jerárquica.

\subsection{Catálogo de competencias y matriz de tributación}
\label{subsec:competencias}

Los procesos de acreditación de carreras de ingeniería requieren que la
institución demuestre, a través de sus propuestas de cátedra, de qué manera
el perfil de egreso se alcanza a lo largo del plan de estudios.
En la práctica, esto implica que cada asignatura debe declarar las
competencias que desarrolla y en qué nivel lo hace, lo que da lugar
a la \emph{matriz de tributación}: una representación que, para cada
competencia del catálogo, muestra en qué asignaturas y con qué grado
de profundidad se trabaja
a lo largo del plan.

El módulo implementa el catálogo de competencias mediante la entidad
\texttt{CompetencyCatalog}, que distingue dos tipos conforme al marco
CONFEDI:

\begin{itemize}
  \item \textbf{Competencias genéricas}: comunes a todas las carreras de
  ingeniería, independientemente de la especialidad, y orientadas al
  ejercicio profesional en sentido amplio.
  \item \textbf{Competencias específicas}: propias de cada disciplina
  de ingeniería y vinculadas
  al dominio técnico particular de la carrera.
\end{itemize}

La relación \texttt{ProposalCompetency} implementa la asociación muchos-a-muchos
entre competencias y propuestas de cátedra, registrando el nivel de
desarrollo declarado por el docente.
Esta información, consolidada a nivel de carrera, constituye la
matriz de tributación que permite verificar:

\begin{enumerate}
  \item Que todas las competencias del perfil de egreso sean trabajadas
  en al menos una asignatura (cobertura completa).
  \item Que la progresión a lo largo del plan sea coherente: las
  competencias deben introducirse en años inferiores y consolidarse
  en los superiores.
  \item Que no existan asignaturas que declaren competencias sin relación
  evidente con su contenido, lo que podría señalar errores de carga
  o inconsistencias en la propuesta.
\end{enumerate}

El endpoint \texttt{GET /competencies/\{id\}/usage} devuelve el conjunto
de propuestas que trabajan una competencia determinada, con el nivel
declarado en cada una, lo que permite al director o a la Comisión
Curricular visualizar la distribución real antes de una visita de
acreditación.
La operación \texttt{POST /competencies/map-plans} ejecuta el mapeo
masivo entre el catálogo de competencias activo y el plan de estudios
vigente, generando o actualizando las asociaciones existentes en una
sola operación en lugar de requerir una vinculación manual asignatura
por asignatura.

%% =============================================================================
\section{Gestión de docentes y carreras}
\label{sec:docentes_carreras}
%% =============================================================================

\subsection{Docentes}
\label{subsec:docentes}

El \textbf{Directivo} es el único rol autorizado para dar de alta a los
docentes en el sistema.
Al registrar a un docente (\texttt{POST /teachers}), se crea
simultáneamente su cuenta de usuario: el correo electrónico
institucional se asigna como nombre de usuario y, por defecto,
también como contraseña inicial.
El sistema exige el cambio de contraseña en el primer inicio de sesión,
de modo que ninguna cuenta quede indefinidamente con credenciales
triviales.

El modelo \texttt{Teacher} almacena nombre, apellido, correo electrónico,
cargo docente y dedicación.
La actualización de estos datos también recae en el Directivo, quien
puede modificarlos en cualquier momento mediante \texttt{PUT /teachers/\{id\}}.
La relación \texttt{ProposalTeacher} vincula docentes a propuestas;
\texttt{TeacherCareer} los vincula a cada carrera con un rol institucional.

El endpoint \texttt{GET /teachers/\{id\}/usage} devuelve las asignaturas
asignadas al docente.
La eliminación está restringida: si el docente tiene propuestas activas,
el backend responde con \texttt{HTTP 409 Conflict}.

El restablecimiento de contraseña puede iniciarse de dos formas:
el propio docente mediante el flujo de recuperación por correo
electrónico (\texttt{POST /auth/forgot-password}),
o el administrador del sistema de manera directa desde el panel
de administración (\texttt{POST /admin/users/\{id\}/reset-password}).

\subsection{Carreras y roles RBAC}
\label{subsec:carreras_rbac}

El modelo \texttt{Career} registra nombre, unidad académica, estado de
acreditación (booleano) y ciclo vigente.
Toda la lógica de propuestas, controles y evidencias está segmentada por
carrera; el rol \texttt{TeacherCareer.role} determina qué puede hacer
cada actor en la plataforma, tal como se resume en la
Tabla~\ref{tab:rbac}:

\begin{table}[ht]
\centering
\caption{Roles y principales permisos por módulo.}
\label{tab:rbac}
\small
\begin{tabular}{lcccc}
\toprule
\textbf{Acción} & \textbf{Directivo} & \textbf{Docente} & \textbf{Secretario} & \textbf{MiembroCC} \\
\midrule
Crear propuesta               & \checkmark &            & \checkmark & \\
Editar propuesta              & \checkmark & \checkmark & \checkmark & \\
Eliminar propuesta              & \checkmark &  & & \\
Habilitar / deshabilitar edición de propuesta              & \checkmark &  & \checkmark & \\
Aprobar propuesta             & \checkmark &            &            & \\
Ejecutar controles rápidos    & \checkmark & \checkmark & \checkmark & \checkmark \\
Ejecutar controles inteligentes & \checkmark & \checkmark & \checkmark & \checkmark \\
Exportar documentación        & \checkmark & \checkmark & \checkmark & \checkmark \\
Gestionar instrumentos        & \checkmark & \checkmark & \checkmark & \checkmark \\
Gestionar Comisión Curricular & \checkmark &            &  & \\
Enviar notificaciones         & \checkmark &            & \checkmark & \checkmark \\
Agregar observaciones         & \checkmark &            & \checkmark & \checkmark \\
\bottomrule
\end{tabular}
\end{table}


%% =============================================================================
\section{Módulo de instrumentos evaluativos}
\label{sec:instrumentos}
%% =============================================================================

Antes de implementar este módulo, el seguimiento de los instrumentos
evaluativos —trabajos prácticos, parciales, finales y recuperatorios—
se realizaba por correo electrónico, sin registro consolidado ni
visibilidad del estado por asignatura.
El módulo centraliza ese seguimiento en la entidad
\texttt{EvaluativeInstrument}, que registra tipo, asignatura,
fecha de entrega, referencia al archivo y estado
(\texttt{pendiente} / \texttt{entregado}).
La Figura~\ref{fig:instrumentos_ciclo} ilustra el ciclo operativo.

\shorthandoff{<>}
\begin{figure}[H]
\centering
\resizebox{0.95\linewidth}{!}{%
\begin{tikzpicture}[
  x=1cm, y=1cm,
  font=\normalsize,
  blk/.style={draw, rounded corners=4pt, minimum width=3.4cm,
               minimum height=1.1cm, align=center,
               fill=gray!10, line width=0.7pt},
  dstep/.style={draw, rounded corners=4pt, minimum width=3.4cm,
                minimum height=1.1cm, align=center,
                fill=Micolor1!12, draw=Micolor1, line width=0.8pt},
  arr/.style={->, >=stealth, line width=1.2pt},
  rol/.style={font=\small\itshape, align=center, text=gray}
]
\node[dstep] (drive) at ( 1.5, 2.5) {Vincular\\carpeta Drive};
\node[blk]  (sube)  at ( 8.5, 2.5) {Subir instrumento/s\\(archivo)};
\node[blk]  (notif) at (15.5, 2.5) {Notificar\\pendientes};

\node[rol] at ( 1.5, 1.2) {Directivo};
\node[rol] at ( 8.5, 1.2) {Docente};
\node[rol] at (15.5, 1.2) {Directivo / Miembro CC};

\draw[arr] (drive) -- (sube);
\draw[arr] (sube)  -- (notif);
\end{tikzpicture}%
}
\caption{Ciclo operativo del módulo de instrumentos evaluativos.
\textit{Fuente: elaboración propia.}}
\label{fig:instrumentos_ciclo}
\end{figure}
\shorthandon{<>}

El Directivo vincula una carpeta de Google Drive a cada asignatura
como paso previo obligatorio para centralizar la entrega de archivos.
Las operaciones de gestión de carpeta son:
\begin{itemize}[noitemsep, topsep=2pt]
  \item \texttt{POST /instruments/folder/link}: asocia una carpeta Drive
    existente a la asignatura;
  \item \texttt{POST /instruments/folder/create}: crea la carpeta
    automáticamente en Drive;
  \item \texttt{DELETE /instruments/folder}: desvincula la asociación.
\end{itemize}
Una vez configurada la carpeta, el Docente sube sus instrumentos mediante
\texttt{POST /instruments/upload}: el archivo queda depositado en la
carpeta Drive de la asignatura y el registro \texttt{EvaluativeInstrument}
queda vinculado tanto a la asignatura como a esa carpeta.

\texttt{GET /instruments} lista los instrumentos con filtros por carrera,
asignatura, tipo y estado.
\texttt{GET /instruments/summary} devuelve la vista agregada para el
director: porcentaje de entrega por asignatura, sin necesidad de consultar
a cada docente por separado.
Cuando quedan instrumentos pendientes, \texttt{POST /instruments/notify}
envía un correo (vía Gmail OAuth) con nombre del docente, asignatura y
plazo, opcionalmente redactado con asistencia LLM.

%% =============================================================================
\section{Módulo de acreditación CONEAU}
\label{sec:acreditacion}
%% =============================================================================

La acreditación ante la CONEAU requiere que la institución presente
evidencia documental organizada según cinco dimensiones: contexto
institucional, plan de estudios, cuerpo académico, alumnos y graduados,
e infraestructura.   Este módulo centraliza esa evidencia y gestiona
el plan de trabajo de la comisión de autoevaluación.

La configuración inicial (\texttt{PUT /accreditation/settings}) establece
el estándar vigente, el ciclo de acreditación, la institución evaluadora
y las fechas del proceso.

\subsection{Registro de evidencias}
\label{subsec:evidencias}

El modelo \texttt{AccreditationEvidenceRegistry} almacena tipo de evidencia,
descripción, dimensión CONEAU, estado y referencia al archivo.
La API expone operaciones \texttt{GET}, \texttt{POST}, \texttt{PATCH} y
\texttt{DELETE} sobre \texttt{/accreditation/evidences}.

El flujo de ingesta masiva opera en dos pasos: primero se escanea una
carpeta de Drive o del filesystem y el modelo de lenguaje propone la
clasificación automática de cada documento por dimensión CONEAU; el usuario
revisa el resultado y confirma.
Las dos variantes del primer paso son:
\begin{itemize}[noitemsep, topsep=2pt]
  \item \texttt{POST /accreditation/ingest-preview}: carpeta de Google Drive;
  \item \texttt{POST /accreditation/ingest-local-preview}: carpeta local.
\end{itemize}
La confirmación se realiza con \texttt{POST /accreditation/ingest},
que crea los registros definitivos en la base de datos.
Este diseño evita incorporar documentos con clasificación incorrecta.

Cada evidencia mantiene:
\begin{itemize}[noitemsep, topsep=2pt]
  \item historial de versiones
    (\texttt{GET /accreditation/evidences/\{id\}/versions});
  \item log de auditoría inmutable
    (\texttt{GET /accreditation/evidences/\{id\}/audit}),
    que registra quién subió, modificó o aprobó el documento y cuándo.
\end{itemize}

\subsection{Plan de trabajo de la autoevaluación}
\label{subsec:plan_trabajo}

Un proceso de autoevaluación institucional no es un evento puntual sino
una secuencia estructurada de responsabilidades distribuidas en el tiempo
entre múltiples actores.
Sin una herramienta que lo gestione, el seguimiento queda fragmentado en
planillas sin visibilidad compartida y el avance de cada actividad solo
se conoce cuando alguien lo pregunta.
Este módulo formaliza esa estructura y la hace observable en tiempo real.

\subsubsection*{Jerarquía del plan de trabajo}

El plan se organiza en cuatro niveles, representados en la
Figura~\ref{fig:workplan_jerarquia}:

\begin{enumerate}
  \item \textbf{Etapa} (\texttt{AccreditationWorkPlanStage}): agrupa fases
    de alto nivel del proceso de autoevaluación
    (\textit{p.\,ej.}\ recolección de evidencias, análisis institucional,
    redacción del informe, visita de pares).
  \item \textbf{Sub-etapa} (\texttt{AccreditationWorkPlanSubStage}): división
    funcional dentro de cada etapa, alineada a las dimensiones CONEAU
    (\textit{p.\,ej.}\ dimensión cuerpo académico, dimensión plan de estudios).
  \item \textbf{Actividad} (\texttt{AccreditationWorkPlanActivity}): unidad
    ejecutable asignada a un responsable principal, con fechas de inicio
    y fin, descripción y lista de colaboradores opcionales.
  \item \textbf{Tarea} (\texttt{AccreditationWorkPlanTask}): ítem atómico
    dentro de una actividad, con su propio responsable, fecha límite
    y estado independiente.
\end{enumerate}

\shorthandoff{<>}
\begin{figure}[H]
\centering
\resizebox{0.50\linewidth}{!}{%
\begin{tikzpicture}[
  x=1cm, y=1cm,
  font=\normalsize,
  lvl/.style={draw, rounded corners=4pt, minimum width=4.0cm,
              minimum height=1.0cm, align=center,
              fill=gray!10, line width=0.7pt},
  sublvl/.style={draw, rounded corners=4pt, minimum width=4.0cm,
                 minimum height=1.0cm, align=center,
                 fill=Micolor1!12, draw=Micolor1, line width=0.8pt},
  arr/.style={->, >=stealth, line width=1.2pt},
  lbl/.style={font=\small\itshape, text=gray}
]
\node[lvl]    (etapa)    at (0, 6.0) {Etapa};
\node[lvl]    (subetapa) at (0, 4.2) {Sub-etapa};
\node[sublvl] (activ)    at (0, 2.4) {Actividad};
\node[sublvl] (tarea)    at (0, 0.6) {Tarea};

\draw[arr] (etapa)    -- (subetapa);
\draw[arr] (subetapa) -- (activ);
\draw[arr] (activ)    -- (tarea);

\node[lbl, right=0.7cm of etapa]    {Responsable, fechas del período};
\node[lbl, right=0.7cm of subetapa] {Dimensión CONEAU, descripción};
\node[lbl, right=0.7cm of activ]    {Responsable, colaboradores, fechas, estado};
\node[lbl, right=0.7cm of tarea]    {Responsable, fecha límite, estado};
\end{tikzpicture}%
}
\caption{Jerarquía del plan de trabajo de autoevaluación.
\textit{Fuente: elaboración propia.}}
\label{fig:workplan_jerarquia}
\end{figure}
\shorthandon{<>}

\subsubsection*{Modelo de estados y sistema de semáforo}

Tanto las actividades como las tareas implementan un ciclo de vida
con los estados:
\texttt{pendiente} → \texttt{en\_progreso} → \texttt{completo},
con la transición automática a \texttt{vencido} cuando la fecha límite
se supera sin que el ítem haya sido marcado como completo.

La interfaz visualiza ese ciclo como un \textbf{semáforo de colores}:

\begin{description}[noitemsep, leftmargin=2.0cm, labelwidth=1.8cm]
  \item[\textcolor{green!60!black}{\textbf{Verde}}]
    Actividad completada dentro del plazo.
  \item[\textcolor{orange!80!black}{\textbf{Naranja}}]
    Actividad en progreso con fecha límite próxima
    (dentro del umbral configurable, por defecto 7 días).
  \item[\textcolor{red!70!black}{\textbf{Rojo}}]
    Actividad vencida: la fecha límite se superó sin completarse.
  \item[\textcolor{gray}{\textbf{Gris}}]
    Actividad pendiente sin fecha próxima.
\end{description}

Este semáforo se calcula en el backend en el endpoint
\texttt{GET /accreditation/workplan} comparando la fecha actual con
\texttt{due\_date} de cada ítem, y es consumido directamente por el
componente de interfaz sin lógica adicional en el frontend.

\subsubsection*{Responsables, colaboradores y acceso}

Cada actividad y cada tarea puede tener asignado un \textbf{responsable
principal} y una lista de \textbf{colaboradores}, ambos tomados del
padrón de docentes existente en el sistema (\texttt{Teacher}).
Esta vinculación no requiere ningún alta adicional: cuando se asigna
un docente como responsable o colaborador de una actividad, ese mismo
usuario obtiene acceso automático a la sección de acreditación
en el frontend, donde solo ve las actividades y tareas que le fueron
asignadas.
El Directivo y el Secretario tienen visión completa del plan; el
Miembro de la Comisión Curricular y el Docente asignado ven
exclusivamente su cartera de ítems.

\subsubsection*{API del módulo}

Las operaciones disponibles sobre el plan de trabajo son:

\begin{itemize}[noitemsep, topsep=2pt]
  \item \texttt{GET /accreditation/workplan}: devuelve el plan completo
    con estados calculados y semáforo para cada ítem;
  \item \texttt{POST /accreditation/workplan/stages}: crea una etapa;
  \item \texttt{POST /accreditation/workplan/substages}: crea una
    sub-etapa dentro de una etapa;
  \item \texttt{POST /accreditation/workplan/activities}: crea una
    actividad con responsable y colaboradores;
  \item \texttt{PATCH /accreditation/workplan/activities/\{id\}}: actualiza
    estado, fechas o responsable de una actividad;
  \item \texttt{POST /accreditation/workplan/activities/\{id\}/tasks}: agrega
    una tarea a una actividad;
  \item \texttt{PATCH /accreditation/workplan/tasks/\{id\}}: actualiza
    estado o fecha límite de una tarea;
  \item \texttt{DELETE /accreditation/workplan/activities/\{id\}}: elimina
    una actividad sin afectar el historial de evidencias asociadas;
  \item \texttt{DELETE /accreditation/workplan/tasks/\{id\}}: elimina
    una tarea sin afectar el historial de evidencias asociadas.
\end{itemize}

Las actividades próximas a vencer o ya vencidas se integran con el módulo
de notificaciones (\S\ref{sec:notificaciones}): el sistema puede generar
automáticamente un correo de alerta con la lista de ítems en riesgo,
redactado con asistencia LLM y enviado por la cuenta Gmail del director.

Cabe destacar que la dimensión \emph{plan de estudios} de la acreditación
no se sustenta únicamente en evidencias documentales: la plataforma
permite construir la matriz de tributación de competencias
(\S\ref{subsec:competencias}) a partir de los datos ya cargados en el
catálogo de competencias y las propuestas de cátedra aprobadas.
Esta matriz, consultable antes de una visita de acreditación, permite
verificar la cobertura del perfil de egreso y detectar brechas entre
lo declarado en el plan de estudios y lo efectivamente desarrollado
en las asignaturas.

%% =============================================================================
\section{Módulo de notificaciones}
\label{sec:notificaciones}
%% =============================================================================

El módulo de notificaciones centraliza el envío de correos institucionales
que antes se redactaban manualmente sin registro ni uniformidad.

\subsection{Integración Gmail OAuth}
\label{subsec:gmail_oauth}

La integración usa el flujo OAuth 2.0 estándar de Google.
Los endpoints involucrados son:
\begin{itemize}[noitemsep, topsep=2pt]
  \item \texttt{GET /notifications/gmail/authorize}: redirige al usuario
    a la pantalla de consentimiento de Google; al completar el flujo,
    el callback guarda el token de acceso y refresco.
  \item \texttt{GET /notifications/gmail/status}: informa si la cuenta
    está conectada.
  \item \texttt{POST /notifications/gmail/disconnect}: revoca el token.
  \item \texttt{POST /notifications/gmail/send}: envía el correo desde
    la cuenta institucional del usuario, no desde un servidor SMTP genérico.
\end{itemize}

\subsection{Asistencia LLM para redacción de correos}
\label{subsec:email_llm}

El endpoint \texttt{POST /notifications/email/ai-assist} recibe el
contexto de la notificación (nombre del docente, asignatura, lista de
campos pendientes o instrumentos no entregados) y genera el cuerpo del
correo en tono institucional formal.
El usuario revisa el texto generado antes de enviarlo, lo que mantiene
el control humano sobre la comunicación.

Los casos de uso cubiertos son: propuestas incompletas fuera de plazo,
instrumentos evaluativos no subidos, actividades del plan de trabajo
próximas a vencer, y resultados de controles normativos con
observaciones que requieren corrección por parte del docente.

%% =============================================================================
\section{Backend — API REST}
\label{sec:backend}
%% =============================================================================

El backend se implementa con \textbf{FastAPI} (Python~3.13.12), un framework
asincrónico que genera automáticamente la documentación interactiva
OpenAPI/Swagger en \texttt{/docs} y valida todos los cuerpos de petición
mediante esquemas \textbf{Pydantic} antes de que lleguen a la lógica de
negocio.
La base de datos es \textbf{SQLite} gestionada con \textbf{SQLAlchemy} en
modo síncrono; las migraciones de esquema se administran con
\textbf{Alembic}, lo que permite evolucionar el modelo de datos sin
pérdida de información entre versiones del sistema.

Todos los endpoints residen en un único \texttt{main.py}, organizados por
recurso mediante prefijos de ruta coherentes con los nombres de los
módulos funcionales.
Esta estructura monolítica simplifica el despliegue en la fase actual del
proyecto, en la que la carga de usuarios y el equipo de desarrollo no
justifican una arquitectura de microservicios; la separación conceptual
entre módulos queda garantizada por la organización interna del código y
por el uso de routers de FastAPI.

\subsection{Autenticación y seguridad}

La autenticación usa \textbf{JWT Bearer token} (algoritmo HS256) con
tiempo de expiración configurable mediante variable de entorno.
Las contraseñas se almacenan exclusivamente como hash \texttt{bcrypt},
nunca en texto plano.
El ciclo completo de gestión de credenciales expone:
\begin{itemize}[noitemsep, topsep=2pt]
  \item \texttt{POST /auth/login}: autenticación y emisión del token;
  \item \texttt{GET  /auth/me}: perfil del usuario autenticado;
  \item \texttt{POST /auth/forgot-password}: solicitud de recuperación
    por correo electrónico;
  \item \texttt{POST /auth/reset-passwords}: reset con token de un solo
    uso enviado por correo;
  \item \texttt{POST /auth/change-my-password}: cambio de contraseña
    para el usuario autenticado, exigido en el primer inicio de sesión.
\end{itemize}

Cada endpoint protegido verifica el token JWT y, adicionalmente, el rol
\texttt{TeacherCareer.role} del usuario para la carrera involucrada.
Esta verificación se realiza en una dependencia reutilizable inyectada
por FastAPI, de modo que la lógica de autorización no se duplica en cada
ruta.

\subsection{Integración con servicios externos}

El backend gestiona tres integraciones con servicios de Google:
\begin{description}[noitemsep, leftmargin=1.6cm, labelwidth=1.4cm]
  \item[Google Docs] Sincronización bidireccional de propuestas de cátedra
    mediante la API de Google Docs; el flujo OAuth se gestiona con
    \texttt{/auth/google/\{authorize,callback\}}.
  \item[Google Drive] Almacenamiento de instrumentos evaluativos en
    carpetas por asignatura; autenticación compartida con Google Docs.
  \item[Gmail] Envío de notificaciones institucionales desde la cuenta del
    director mediante OAuth~2.0 (\S\ref{sec:notificaciones}).
\end{description}

La clave \texttt{OPENAI\_API\_KEY} habilita las funciones de generación,
reformulación y evaluación mediante LLM (controles inteligentes, asistencia
de redacción, clasificación de evidencias CONEAU).
Todas las credenciales se leen desde variables de entorno; el archivo
\texttt{.env} y la base de datos están cubiertos por \texttt{.gitignore}.

\subsection{Tareas en segundo plano}

Las operaciones de larga duración —importación de DOCX, ejecución de
controles inteligentes, ingesta masiva de evidencias— se ejecutan como
\texttt{BackgroundTasks} de FastAPI.
Esto permite que el endpoint devuelva \texttt{HTTP~202 Accepted}
inmediatamente mientras el procesamiento continúa en segundo plano,
sin bloquear al cliente y sin necesidad de un broker de colas externo
en la fase actual del proyecto.

\subsection{Resumen de endpoints por módulo}

La Tabla~\ref{tab:endpoints} resume los grupos de endpoints con la
cantidad real de rutas en cada módulo.

\begin{longtable}{p{4cm}rp{10cm}}
\caption{Grupos de endpoints de la API REST.}
\label{tab:endpoints}\\
\toprule
\textbf{Módulo} & \textbf{N.\textsuperscript{o}} & \textbf{Rutas representativas} \\
\midrule
\endfirsthead
\multicolumn{3}{l}{\small\textit{(continuación)}} \\
\toprule
\textbf{Módulo} & \textbf{N.\textsuperscript{o}} & \textbf{Rutas representativas} \\
\midrule
\endhead
\midrule
\multicolumn{3}{r}{\small\textit{continúa en la página siguiente}} \\
\endfoot
\bottomrule
\endlastfoot
Propuestas         & 17 & \texttt{GET/POST /proposals},
                          \texttt{PATCH /proposals/\{id\}} \\
Google Workspace   & 13 & \texttt{/auth/google/authorize},
                          \texttt{/proposals/\{id\}/sync-gdoc} \\
Plan de estudios   & 12 & \texttt{/study-plans}\\
Acreditación       & 14 & \texttt{/accreditation/evidences},
                          \texttt{/accreditation/ingest} \\
Instrumentos       & 10 & \texttt{/instruments},
                          \texttt{/instruments/upload} \\
Controles intelig. &  8 & \texttt{/intelligent-controls},
                          \texttt{/proposals/\{id\}/intelligent-controls/run} \\
Notificaciones     &  6 & \texttt{/notifications/gmail/send},
                          \texttt{/notifications/email/ai-assist} \\
Competencias       &  5 & \texttt{/competencies},
                          \texttt{/competencies/\{id\}/usage} \\
Docentes           &  5 & \texttt{/teachers},
                          \texttt{/teachers/\{id\}/usage} \\
Carreras           &  6 & \texttt{/careers},
                          \texttt{/careers/\{id\}/committee} \\
Autenticación      &  5 & \texttt{/auth/login},
                          \texttt{/auth/forgot-password} \\
IA generativa      &  4 & \texttt{/suggest},
                          \texttt{/ai-generate},
                          \texttt{/ai-reformulate} \\
Búsqueda semántica &  1 & \texttt{POST /search} \\
\midrule
\textbf{Total} & \textbf{113} & \\
\end{longtable}

%% =============================================================================
\section{Frontend — Interfaz de usuario}
\label{sec:frontend}
%% =============================================================================

El frontend es una SPA (Single Page Application) construida con React~18 y Vite~7.
El enrutamiento del lado del cliente usa \texttt{react-router-dom};
las llamadas a la API se realizan con \texttt{Axios}.
La visibilidad de cada sección y las acciones disponibles dentro de ella
se determinan en tiempo de ejecución según el rol del usuario autenticado.

\subsection*{Secciones y acceso por rol}

La Tabla~\ref{tab:frontend_roles} resume las secciones de la plataforma
y los roles que pueden acceder a cada una.

\begin{longtable}{p{5cm}cccc}
\caption{Secciones del frontend y acceso por rol.}
\label{tab:frontend_roles}\\
\toprule
\textbf{Sección} & \textbf{Directivo} & \textbf{Docente} & \textbf{Secretario} & \textbf{MiembroCC} \\
\midrule
\endfirsthead
\multicolumn{5}{l}{\small\textit{(continuación)}} \\
\toprule
\textbf{Sección} & \textbf{Directivo} & \textbf{Docente} & \textbf{Secretario} & \textbf{MiembroCC} \\
\midrule
\endhead
\bottomrule
\endlastfoot
Dashboard (inicio)           & \checkmark &  & \checkmark & \checkmark \\
Propuestas                   & \checkmark & \checkmark & \checkmark & \checkmark \\
\quad Importar propuesta      & \checkmark &            & \checkmark & \\
\quad Editar propuesta        & \checkmark & \checkmark & \checkmark & \\
\quad Configurar Google Drive & \checkmark &            & \checkmark & \\
Control de propuestas        & \checkmark & \checkmark & \checkmark & \checkmark \\
\quad Control rápido          & \checkmark & \checkmark & \checkmark & \checkmark \\
\quad Control inteligente     & \checkmark & \checkmark & \checkmark & \checkmark \\
\quad Config. ctrl. inteligente & \checkmark &          & \checkmark & \\
Plan de estudios             & \checkmark &            & \checkmark & \checkmark \\
Instrumentos evaluativos     & \checkmark & \checkmark & \checkmark & \checkmark \\
Docentes                     & \checkmark &            & \checkmark & \checkmark \\
\quad Crear / modificar docente & \checkmark &          &            & \\
\quad Eliminar docente        & \checkmark &            &            & \\
\quad Enviar notificaciones   & \checkmark &            & \checkmark & \checkmark \\
Acreditación                 & \checkmark &            & \checkmark & (*) \\
Búsqueda semántica           & \checkmark &            & \checkmark & \checkmark \\
\end{longtable}

\noindent\small(*) Los docentes asignados como responsables o colaboradores
en el plan de trabajo de acreditación acceden únicamente a sus actividades
y tareas vinculadas (\S\ref{subsec:plan_trabajo}).
\normalsize

\subsection*{Descripción de las secciones principales}

\begin{description}[noitemsep, leftmargin=0pt]
  \item[\textbf{Dashboard}]
    Pantalla de inicio común a todos los roles menos al de Docente.
    Muestra el resumen de propuestas (aprobadas, pendientes, con
    observaciones), Docente (con información completa o incompleta) y Asignaturas en el Plan de Estudios..

  \item[\textbf{Propuestas}]
    Listado de propuestas con filtros por carrera, estado y docente.
    El Directivo y el Secretario pueden importar desde DOCX o Google Docs
    y configurar la carpeta Drive asociada.
    El Docente accede únicamente a sus propias propuestas y puede editarlas
    campo a campo con asistencia LLM (\textit{Sugerir}, \textit{Generar},
    \textit{Reformular}, o \textit{Escribir completamente}).

  \item[\textbf{Control de propuestas}]
    Ejecuta controles rápidos (validaciones normativas fijas) o controles
    inteligentes (evaluación por LLM) sobre una propuesta.
    El reporte muestra un semáforo por control con estado, justificación
    textual y sugerencia de mejora.
    La configuración de los controles inteligentes —prompts, modelos,
    modo de ejecución— es exclusiva del Directivo y el Secretario.

  \item[\textbf{Plan de estudios}]
    Árbol jerárquico editable: año $\to$ cuatrimestre $\to$ asignatura,
    con sus cargas horarias y correlatividades.
    Incluye la visualización de la matriz de tributación de competencias
    (\S\ref{subsec:competencias}).
    Accesible para Directivo, Secretario y Miembro de la CC.

  \item[\textbf{Instrumentos evaluativos}]
    Visible para todos los roles con distintos alcances.
    El Directivo vincula carpetas de Google Drive por asignatura.
    El Docente sube y gestiona sus propios instrumentos.
    El Directivo, el Secretario y los Miembros de la CC acceden al resumen agregado de entregas
    y pueden enviar notificaciones a docentes con instrumentos pendientes.

  \item[\textbf{Docentes}]
    Gestión del padrón institucional.
    El Directivo crea y modifica docentes; es el único que puede
    eliminarlos.
    El Secretario y el Miembro de la CC pueden consultar el padrón y
    enviar notificaciones pero no modificar datos personales.

  \item[\textbf{Acreditación}]
    Registro de evidencias por dimensión CONEAU y plan de trabajo de la
    autoevaluación con semáforo de estados.
    El Directivo y el Secretario tienen visión y edición completa.
    Los docentes designados como responsables o colaboradores de
    actividades específicas acceden únicamente a su cartera de ítems.

  \item[\textbf{Búsqueda semántica}]
    Campo de consulta en lenguaje natural sobre el corpus institucional
    indexado (propuestas, normativa, competencias, evidencias).
    Devuelve los fragmentos más relevantes con \emph{score} de similitud
    y permite inyectarlos como contexto al asistente LLM.
    Disponible para Directivo, Secretario y Miembro de la CC.
\end{description}

%% =============================================================================
\section{Seguridad, trazabilidad y auditoría}
\label{sec:seguridad}
%% =============================================================================

\subsection{Autenticación y control de acceso}
\label{subsec:seg_auth}

Todos los endpoints protegidos requieren un JWT Bearer token (HS256)
cuya expiración es configurable.
Las contraseñas se almacenan exclusivamente como hash \texttt{bcrypt},
nunca en texto plano.
La recuperación de contraseña emite un token de un solo uso enviado
por correo; el token expira tras su primer uso.
El control de acceso granular se implementa mediante
\texttt{TeacherCareer.role}: cada endpoint verifica que el usuario
autenticado posea el rol requerido antes de ejecutar la operación.

\subsection{Protección de credenciales y configuración}
\label{subsec:seg_creds}

Todas las credenciales sensibles:

(\texttt{OPENAI\_API\_KEY},
\texttt{GOOGLE\_CLIENT\_ID}, \texttt{GOOGLE\_CLIENT\_SECRET},
\texttt{JWT\_SECRET}) se leen exclusivamente desde variables de entorno
cargadas con \texttt{python-dotenv}.
El archivo \texttt{.env} y la base de datos \texttt{database.db}
están cubiertos por \texttt{.gitignore} para evitar su inclusión
en el repositorio.

\subsection{Trazabilidad y auditoría}
\label{subsec:seg_auditoria}

El sistema implementa dos niveles de auditoría:

\begin{enumerate}
  \item \textbf{Log de evidencias}: \texttt{AccreditationEvidenceAuditLog}
registra cada operación sobre un documento de acreditación con
usuario, acción y marca temporal.
  \item \textbf{Log de acciones de agentes MCP}: el Servidor de Auditoría
(§\ref{subsec:mcp_auditoria}) escribe un registro inmutable por cada
herramienta invocada, incluyendo \texttt{hash} del payload para detectar
alteraciones posteriores.
\end{enumerate}

Las propuestas aprobadas pueden bloquearse contra edición posterior
mediante \texttt{PATCH /proposals/\{id\}/lock} o en bloque al cierre
del ciclo con \texttt{POST /proposals/bulk-lock}.

\subsection{Validación de entradas y mitigación de inyección}
\label{subsec:seg_inputs}

Todos los cuerpos de petición son validados automáticamente por esquemas
Pydantic antes de llegar a la lógica de negocio.
Las URLs de Google Docs son sanitizadas mediante expresión regular para
extraer el \texttt{doc\_id} antes de llamar a la API de Google,
evitando SSRF por URLs maliciosas.
En el contexto del agente LLM, el contenido controlado por el usuario
(texto de propuestas) nunca se concatena directamente en el system prompt;
se inyecta en la sección de usuario del mensaje para preservar el
aislamiento de instrucciones y mitigar ataques de \emph{prompt injection}.
Las herramientas disponibles para el modelo están definidas estáticamente
en cada servidor MCP; el agente no puede invocar herramientas fuera
de ese conjunto.

%% =============================================================================
\section{Decisiones de diseño y trade-offs}
\label{sec:decisiones_disenio}
%% =============================================================================

Todo sistema de software implica decisiones que no tienen una respuesta
universalmente correcta: la elección entre dos alternativas técnicas
igualmente válidas depende del contexto, los recursos disponibles y el
horizonte de evolución del proyecto.
Esta sección documenta las decisiones no triviales tomadas durante el
desarrollo de \textsc{macau}, las alternativas consideradas y la
justificación de cada elección.
Aquellas decisiones que se identifican como candidatas a revisión
en una fase de escala se señalan explícitamente.

\subsection*{Base de datos: SQLite}

\textbf{Alternativas consideradas:} PostgreSQL, MySQL.

Se eligió \textbf{SQLite} por su modelo de despliegue sin servidor: la
base de datos es un único archivo en disco que no requiere instancia
separada, configuración de red ni gestión de usuarios de base de datos.
En el contexto actual —una institución con un único plan de estudios
activo por carrera, menos de diez carreras y un volumen de propuestas que no supera
las pocas centenas de filas— el overhead operativo de PostgreSQL no
se justifica.
SQLAlchemy abstrae la capa de acceso a datos con una interfaz uniforme,
lo que permite migrar a PostgreSQL sin cambios en la lógica de negocio
si el volumen de usuarios o la necesidad de concurrencia de escritura lo
requieren en el futuro.

\subsection*{Índice vectorial: hnswlib}

\textbf{Alternativas consideradas:} ChromaDB, Weaviate, Pinecone.

El módulo RAG utiliza \textbf{hnswlib} como motor de búsqueda aproximada
de vecinos más cercanos.
A diferencia de ChromaDB o Weaviate, hnswlib no requiere un proceso
servidor independiente: el índice se persiste en disco como un archivo
binario y se carga en memoria al iniciar el backend.
Esta decisión elimina una dependencia de infraestructura y simplifica
el despliegue en entornos sin contenedores ni orquestación.
La contrapartida es que el índice no soporta actualizaciones parciales
eficientes a gran escala; sin embargo, en el corpus institucional actual
(propuestas, normativa, competencias y evidencias) el tamaño no supera
los miles de fragmentos, rango en el que hnswlib opera con latencia
de consulta por debajo de los 10~ms.

\subsection*{Modelo de lenguaje: API de OpenAI con modelo configurable}

\textbf{Alternativas consideradas:} Llama~3, Mistral, modelos locales
vía Ollama.

Se optó por la \textbf{API de OpenAI} por tres razones convergentes:
calidad de razonamiento en español suficiente para evaluación académica
estructurada, estabilidad de la API con versionado explícito, y ausencia
de requerimiento de GPU local.
A ello se suma que el autor contaba con créditos preexistentes en la
plataforma, lo que eliminó la barrera de adopción durante el desarrollo
del prototipo.

El sistema no fija un único modelo para todas sus funciones, lo que
permite ajustar la relación costo/calidad según la tarea.

Las funcionalidades de \textbf{asistencia de redacción} —sugerencias de
mejora, generación de texto desde cero y reformulación de resultados de
aprendizaje— utilizan un modelo de alta capacidad del catálogo de OpenAI,
configurado al momento del despliegue, cuyo rendimiento en generación de
texto libre en español resulta adecuado para el contexto académico.

El módulo de \textbf{controles inteligentes} delega la elección del modelo
al Directivo en cada ejecución, conforme a lo descrito en
\S\ref{subsec:controles_inteligentes}.

Los modelos de código abierto ejecutados localmente representan una
alternativa válida para despliegues con restricciones de privacidad
o sin conectividad.
La arquitectura MCP permite sustituir el proveedor LLM sin modificar
la lógica de los agentes, lo que reduce el costo de una migración futura.

\subsection*{Modelo de embeddings: text-embedding-3-small}

\textbf{Alternativas consideradas:} \texttt{text-embedding-ada-002},
modelos multilinguales (e5-multilingual, bge-m3).

El modelo \texttt{text-embedding-3-small} de OpenAI ofrece una relación
calidad/costo superior a su predecesor \texttt{ada-002} con la mitad del
costo por token y dimensionalidad configurable.
Comparte el mismo proveedor que el modelo de chat, lo que simplifica la
gestión de claves y facturación.
Los modelos multilinguales de código abierto son una alternativa para
escenarios sin conectividad, pero requieren infraestructura de cómputo
adicional y su rendimiento en corpus institucionales en español
específico es variable.

\subsection*{Arquitectura del backend: monolito}

\textbf{Alternativas consideradas:} microservicios, arquitectura hexagonal.

El backend concentra todos los endpoints en un único módulo
\texttt{main.py}.
Esta decisión responde al estadio actual del proyecto: un equipo de
desarrollo unipersonal en fase de prototipo necesita iterar rápido,
y la separación en microservicios introduce overhead de coordinación que
no está justificado mientras el sistema corre en un único servidor.
La separación lógica por módulo se mantiene mediante prefijos de ruta
y routers de FastAPI, lo que facilita una futura extracción a
microservicios cuando el volumen de usuarios o la necesidad de
despliegue independiente por módulo lo justifiquen.

\subsection*{Protocolo de agentes: MCP}

\textbf{Alternativas consideradas:} llamadas directas a funciones Python,
LangChain, LlamaIndex.

El \emph{Model Context Protocol} se eligió sobre llamadas directas porque
introduce una capa de trazabilidad nativa: toda invocación de herramienta
queda registrada con el hash del payload, el agente invocante y la marca
temporal, sin instrumentación adicional del código.
Frente a frameworks como LangChain, MCP define un contrato de transporte
(JSON-RPC 2.0 por stdio o HTTP+SSE) que es independiente del modelo y
del lenguaje de implementación, lo que otorga interoperabilidad futura
con otros clientes MCP y reduce el acoplamiento con una biblioteca
específica del ecosistema Python.

\subsection*{Autenticación: JWT local}

\textbf{Alternativas consideradas:} Auth0, Supabase Auth, Keycloak.

Se implementó autenticación \textbf{JWT propia} (HS256) en lugar de
delegar en un proveedor de identidad externo.
La decisión responde a dos criterios: el contexto institucional no tiene
SSO existente con el que integrarse, y los proveedores de identidad
gestionados (Auth0, Supabase) introducen una dependencia de pago y
conectividad externa que no es deseable en una herramienta de uso
interno.
El flujo implementado cubre los requisitos de seguridad del escenario:
tokens con expiración, hash bcrypt de contraseñas, recuperación por
correo con token de un solo uso y cambio de contraseña obligatorio en
el primer inicio de sesión.
Si la institución adoptara en el futuro un IdP corporativo (Azure AD,
Google Workspace), la dependencia \texttt{get\_current\_user} de FastAPI
puede reemplazarse sin modificar la lógica de negocio.

\subsection*{Integración de correo: Gmail OAuth}

\textbf{Alternativas consideradas:} SMTP genérico, SendGrid, Mailgun.

El envío de notificaciones usa la \textbf{API de Gmail} mediante OAuth~2.0
en lugar de un servidor SMTP externo.
Esto garantiza que los correos se envíen desde la cuenta institucional
real del Director, Secretario o Miembro de la CC, con el nombre del remitente y el historial en su
buzón de enviados, lo que otorga confianza y trazabilidad que un
relay SMTP genérico no puede proveer.
La contrapartida es la fricción inicial del flujo de consentimiento
OAuth; sin embargo, el token de refresco se almacena en el backend y
se renueva automáticamente, de modo que el Director, Secretario o Miembro de la CC solo debe autorizar
una vez.

\subsection*{Resumen de decisiones y deuda técnica}

La Tabla~\ref{tab:decisiones} consolida las decisiones descritas,
indicando cuáles se recomiendan revisar ante un escenario de escala.

\begin{longtable}{p{4.5cm}p{4.5cm}p{6cm}}
\caption{Decisiones de diseño y candidatos a revisión futura.}
\label{tab:decisiones}\\
\toprule
\textbf{Decisión} & \textbf{Opción elegida} & \textbf{Candidato a revisión} \\
\midrule
\endfirsthead
\multicolumn{3}{l}{\small\textit{(continuación)}} \\
\toprule
\textbf{Decisión} & \textbf{Opción elegida} & \textbf{Candidato a revisión} \\
\midrule
\endhead
\bottomrule
\endlastfoot
Base de datos        & SQLite             & Migrar a PostgreSQL al superar múltiples instituciones o concurrencia alta \\
Índice vectorial     & hnswlib            & ChromaDB/Weaviate si el corpus supera decenas de miles de fragmentos \\
Modelo LLM           & API OpenAI (modelo configurable por módulo) & Modelo local si se requiere privacidad estricta o sin conectividad \\
Embeddings           & text-embedding-3-small & Modelo local multilingüe en despliegues sin API externa \\
Arquitectura backend & Monolito           & Microservicios si se requiere escala independiente por módulo \\
Protocolo agentes    & MCP                & Estable; escala con el protocolo sin cambios de arquitectura \\
Autenticación        & JWT local          & IdP corporativo si la institución adopta SSO \\
Integración correo   & Gmail OAuth        & SMTP corporativo si se exige independencia de Google \\
\end{longtable}
